 \documentclass{oqireport}

%% Language and font encodings
\usepackage[english]{babel}
\usepackage[T1]{fontenc}
\usepackage{array}

%% Sets page size and margins
\usepackage[a4paper,top=3cm,bottom=2cm,left=3cm,right=3cm,marginparwidth=1.75cm]{geometry}

%% Useful packages
\usepackage{amsmath}
\usepackage{graphicx}
\usepackage[colorinlistoftodos]{todonotes}
\usepackage[colorlinks=true, allcolors=blue]{hyperref}
\usepackage{algpseudocode}
\usepackage{algorithm}
\usepackage{ulem}
\usepackage{multicol}
\usepackage{booktabs}
\usepackage{multirow}
\usepackage{tikz}
\usetikzlibrary{shadows,arrows.meta,positioning,calc,decorations.pathreplacing}
\usepackage{adjustbox}
\usepackage{makecell}

\title{Food Production Optimization}
\subtitle{Full Proposal}
\author{Edoardo Spigarolo}


\begin{document}

% Instructions to be deleted before submission
\ \\
\ \\
\ \\
\textit{
Instructions and background: Another test for suggested changes in the text.
\begin{itemize}
    \item For each section the envisaged content is specified in italic. All of these comments should be removed before submitting the final document.
    \item Please be concise and respect the strict page limit of 15 pages for the main part of the full proposal (excluding the Methods section, the Team Presentation section and References).
    \item The purpose of the Full Proposal is to present a compelling story for your use case. Based on this a decision regarding the feasibility of a meaningful implementation on a quantum simulator with further OQI support in Phase 3 will be made.\\
    The addressed problem and its broader societal context are first described, and then linked to an isolated computational problem that is formulated mathematically. An analysis of existing (classical) computational approaches is provided, highlighting their limitations and identifying specific bottlenecks in detail. A quantum computing solution is then proposed that has the potential to overcome these challenges. A comprehensive literature review supports all of these elements. Crucially, the proposal includes a clear and well-founded justification for employing quantum computing in the given context.
    Subsequently, a proof of concept is outlined, including a robust benchmarking strategy and a detailed resource estimation.
    Finally, the potential impact of the proposed quantum computing approach is analysed in a broader context.
\end{itemize}
}

\newpage

\maketitle

\begin{abstract}

 Nutrition and the environment are inextricably linked. Climate change adversely affects yields and nutritional quality of major staple crops, raising the risk of food insecurity and malnutrition. Conversely, our food systems drive an estimated 80\% of deforestation and a third of greenhouse gas emissions. As populations and incomes grow in developing countries, demand for animal-source foods will rise, placing further pressure on meeting climate and biodiversity loss goals. \cite{Fanzo2022, UNFSS_Planet_2023, Sylvester2024}

With the rapid progress in quantum computing, various optimization tasks in industry, energy, and other societal applications may take advantage of these novel methodological opportunities \cite{ajagekar_quantum_2019,abbas_quantum_2023,osaba_hybrid_2024}.

To this purpose, a set of well established ways to convert Mixed Integer Linear Programming (MILP) or Mixed Integer Quadratic Programming (MIQP) optimization problems into Quadratic Binary Unconstrained Optimization (QUBO) problems -- those that are naturally handled by quantum computers and annealers -- is now available \cite{naghmouchi_mixed_2024,franco_efficient_2023,karimi_practical_2019}.

Our goal is to demonstrate the usage of such methods in the resolution of a food optimization problem, showing how they constitute a valid alternative to classical solvers in predefined instances.\cite{karimi_practical_2019,ajagekar_quantum_2020,ajagekar_quantum_2021,chang_hybrid_2022,wang_quantum-inspired_2022,yonaga_quantum_2022,franco_efficient_2023,paterakis_hybrid_2023,ellinas_hybrid_2024,naghmouchi_mixed_2024}.
\end{abstract}

\subsubsection*{Relevant SDGs}
\begin{itemize}
    \item \textbf{SDG 2: Zero Hunger}
    \item SDG 3: Good Health and Well-being
    \item SDG 12: Responsible Consumption and Production
    \item SDG 13: Climate Action
\end{itemize}

\section{Description of the problem and its societal context}

The global food system faces mounting pressure to provide nutritious, affordable, and environmentally sustainable food. According to the Food and Agriculture Organization (FAO), global hunger has risen to 9.2\% of the world population in 2022, compared with 7.9\% in 2019 \cite{FAO_hunger}. Conflict and insecurity, economic shocks, and extreme weather events are the main drivers of acute food insecurity. Current food systems are characterized by significant limitations that contribute to the triple burden of malnutrition: undernutrition, micronutrient deficiencies, and overweight/obesity. Research indicates that 88\% of countries experience the coexistence of multiple burdens of malnutrition, often associated with diets high in saturated fat, sugar, and processed foods while being low in fiber, fruits, and vegetables \cite{ZOU2022100624}.

Food systems are complex networks of activities, actors, resources, and environments encompassing the production, processing, distribution, consumption, and disposal of food products. Comprehending these networks is arduous as they are constantly evolving, and are interconnected with broader social, economic and policy systems. To achieve Sustainable Development Goal (SDG) 2 \cite{UN_SDG2}, the FAO affirms that a better understanding of food systems requires a local-to-global perspective.

Climate change has a very strong impact on food systems, exacerbating hunger and malnutrition issues, especially in regions that are already vulnerable. It triggers extreme weather events, the emergence and spread of pests and diseases, and it compromises the adaptation of traditionally grown foods, resulting in reduced crop yields \cite{IPCC_food_security}. Current projections indicate a necessary increase of food production per hectare by almost 60 percent by 2050 to meet the needs of the projected global population of 10 billion—assuming a static agricultural land \cite{GESDA_radar}. In this context, increasing productivity of agricultural land is not sufficient to address the core human need for food. Even now, in countries with low and middle incomes, diets are known to be lacking in micro-nutrients including iron, zinc, folate, vitamin A, calcium, and vitamin B12. As populations continue to grow, innovative solutions are much needed to devise sustainable agricultural practices that provide for more nutritious diets, while respecting planetary boundaries.

On the other hand, food systems also contribute significantly to climate change, accounting for 34\% of global greenhouse gas (GHG) emissions. While the production of food is a large contributor, greenhouse gas emissions associated with food transport also keep increasing \cite{Pradhan2022}. Both the UNFCCC and the FAO are working not only towards climate-resilient and low-emission agricultural practices, but also acknowledge the considerable potential to maximize co-benefits of adaptation and mitigation of global food supply chains \cite{FAO_climate_strategy}.

Recent large-scale disruptions like the Panama Canal drought have reinforced awareness of the vulnerability of the global food supply chains \cite{Nelson2021}, in critical chokepoints represented by maritime corridors such as straits and canals, and coastal and inland transport infrastructures in major crop importing-exporting regions. Accurate modelling of food systems will help provide the basis for improving our ability to produce nutritious food sustainably and foster resilience and respond to unforeseeable food systems disruptions.

Key challenges include:
\begin{itemize}
    \item Ensuring food security and nutritional adequacy for growing populations
    \item Reducing the environmental footprint of agriculture
    \item Balancing economic viability for producers with affordability for consumers
    \item Promoting diversity in crop production to enhance resilience and health outcomes
    \item Addressing the fundamental issue that healthy diets remain unaffordable for many populations
    \item Improving the resilience of global food supply chains against potential disruptions
\end{itemize}

\textbf{Who is affected:} Both local and global populations are impacted, with particular vulnerability among low-income communities and regions with limited arable land. 

\textbf{Geographies and demographics:} The affected populations are widespread, but the solution depends on the chosen region, types of crops, and local producer profiles. By 2050, 68\% of the world's population will live in urban regions, consuming 80\% of total food production while cities occupy only 3\% of Earth's surface \cite{ZOU2022100624}. Climate change will further exacerbate these challenges, with 70\% of studies indicating crop yield declines by the 2030s \cite{ZOU2022100624}.

\textbf{Criticality and risks:} Risks include persistent malnutrition, unsustainable land use, and economic disparities. When we consider crop planning that is coordinated across the entire supply chain, we face a fundamental optimization challenge: given a fixed amount of farms and arable land, how can we best allocate these resources to feed the maximum number of people in a nutritionally intelligent way, at the lowest cost to the population, while minimizing environmental harm?


\section{Description of the computational challenges}

\subsection{Mathematical formulation of the problem}
We address the problem of centralized planning of sustainable food production in a pre-established area. To this end, we take a smaller instance of the problem presented in \cite{esteso_sustainable_2023} and we model it as a multi-objective optimization problem.
The subset contains the following objectives:
\begin{itemize}
    \item maximize nutritional value
    \item maximize nutrient density
    \item maximize affordability
    \item maximize sustainability
    \item minimize environmental impact
\end{itemize}
Our objective can in a first approximation be expressed as:
\begin{align}
    &Z = w_1G_1 + w_2 G_2 + \dots \\
    \text{subject to } &\text{Area, Crops}
    \label{eq:generalized objective}
\end{align}
where $G_i$ are the goals above and $w_i$ are the weights we associate to them.

To solve this multi-objective formulation, we transform it into a Mixed-Integer Linear Program (MILP) \cite{Dantzig1951}, a well-studied class of optimization problems solved efficiently by state-of-the-art solvers like Gurobi \cite{gurobi2023}, as demonstrated in \cite{esteso_sustainable_2023}.

Each goal $G_i$ is formulated as a linear combination of decision variables, weighted according to stakeholder preferences or policy priorities. The constraints encode physical limitations (available land area), agronomic requirements (minimum planting areas, crop compatibility), nutritional targets (food group diversity, dietary adequacy), and environmental bounds (emissions thresholds, resource limits). The parameters that define these constraints -- including nutritional scores, environmental impacts, land availability, and dietary requirements -- are derived from established databases and nutritional guidelines, ensuring the model reflects real-world agricultural and nutritional considerations.

\paragraph{Decision Variables}
Let $A_{f,c}$ denote the area in hectares assigned to food $c$ on farm $f$, whether it is a crop or an animal.
Introduce auxiliary binary variables $Y_{f,c}$, where:
\begin{equation}
Y_{f,c} = \begin{cases}
1 & \text{if food $c$ is planted at farm $f$,}\\
0 & \text{otherwise.}
\end{cases}
\end{equation}
Additionally, introduce global binary variables $U_c$ to track unique food selection across all farms:
\begin{equation}
U_c = \begin{cases}
1 & \text{if food $c$ is planted on at least one farm,}\\
0 & \text{otherwise.}
\end{cases}
\end{equation}
The variables $U_c$ are essential for correctly enforcing food group diversity constraints, which require counting the number of \textit{distinct} foods selected rather than the total number of plot-crop assignments.
With this formulation the problem is now expressed as a MILP: the area $A_{f,c} \in [0, A_{max}]$\footnote{The minimum area here is set to 0 as it is the lowest bound when not planted, but the interval could also be expressed as the disjoint $\{0\}\cup[A_{c, min}, A_{max}]$ if the minimum area is relevant, for example when considering to build shelter for animals.} can then be replaced by the combination of a (semi-)continuous variable $A_{f,c} \in [-\infty, A_{max}]$ and the binary variable $Y_{f,c}$ as the product $Y_{f,c} \cdot A_{f,c} \in [0, A_{max}]$. 

\paragraph{Parameters}
To provide a model that is as close as possible to real-world scenarios, the following parameters will be used to characterize our data:
\begin{itemize}
    \item Nutritional value score for food $c$ (higher is better).
    \item Nutrient density score for food $c$ (higher is better).
    \item Environmental impact score for food $c$ (lower is better).
    \item Affordability score for food $c$ (higher is better).
    \item Total land available at farm $f$.
    \item Minimum area to be planted for food $c$ (if selected).
    \item Set of foods belonging to food group $g$.
    \item Minimum and maximum number of foods to produce per food group $g$, following from dietary constraints.
\end{itemize}



\subsection{Classical computational approaches \cite{CLAUTIAUX2025707}}
\label{sub:classical}

\input{diagram_hardness}


Classical MILP solvers have achieved remarkable sophistication through the systematic integration of several foundational algorithmic components, such as branch-and-bound methodology, cutting-plane techniques (branch-and-cut), column generation (branch-and-price), primal heuristic strategies \cite{nemhauser1988integer, nemhauser1999integer}. The branch-and-bound algorithm forms the cornerstone of modern MILP solvers, employing a recursive partitioning strategy that systematically explores the integer feasible region through variable fixation and continuous LP relaxations at each node of the enumeration tree \cite{achterberg2007constraint, wolsey1998integer}. The algorithm's efficiency derives from its pruning mechanism, whereby nodes with relaxation bounds that cannot improve upon the current incumbent solution are eliminated, thereby substantially reducing the computational search space.


Leading commercial implementations, include IBM CPLEX and Gurobi Optimizer \cite{cplex2023, gurobi2023}.

Despite their remarkable sophistication, classical MILP solvers face several fundamental bottlenecks that can significantly limit their performance and scalability across diverse problem classes:

\begin{itemize}


\item LP relaxation tightness represents a primary concern, as weak relaxations with large integrality gaps force the branch-and-bound algorithm to explore extensive enumeration trees before effective pruning can occur~\cite{nemhauser1999integer}.

\item While cutting-plane methods can strengthen these relaxations, the computational overhead of cut generation -- including the solution of auxiliary linear programs and execution of combinatorial separation algorithms -- can itself become prohibitively expensive, particularly when yielding diminishing marginal improvements~\cite{mitchell2002cutting}.

\item Branch-and-bound tree explosion constitutes another critical limitation, where suboptimal variable selection and node selection heuristics may cause exponential growth in the number of subproblems, consuming substantial computational memory and potentially exhausting available system resources before optimal solutions are identified~\cite{linderoth2005anatomy}.

\item In column generation frameworks (branch-and-price), the decomposition of large-scale MILPs into master problems and pricing subproblems introduces iterative solution procedures that can themselves constitute computationally intensive linear or mixed-integer programs, with convergence potentially stalling when pricing subproblems prove difficult or generate only marginally beneficial columns~\cite{lubbecke2005selected}.

\item Primal heuristic strategies, while effective at rapidly identifying feasible solutions, may become trapped in suboptimal regions of the solution space, and the development of problem-specific heuristic procedures requires substantial domain expertise that cannot always be automated effectively~\cite{berthold2013heuristic}.

\item Numerical stability issues arising from ill-conditioned constraint matrices necessitate frequent basis refactorizations in LP solution procedures, while aggressive preprocessing strategies may inadvertently eliminate structural properties essential for specialized cutting-plane generation or decomposition approaches~\cite{gamrath2015structure}.

\item Finally, parallel scalability remains constrained by synchronization overhead and load-balancing challenges across multiple computational threads, particularly in the presence of intensive cut generation or complex heuristic procedures~\cite{shinano2018parallelization}. 

\end{itemize}

\subsection{Description of the potential quantum computing solution}\label{PoC}

The optimization model of interest (e.g.\ \cite{esteso_sustainable_2023}) is posed as a mixed-integer linear program (MILP): a mathematical program that minimizes a linear objective function subject to linear equality and inequality constraints, where some decision variables are restricted to integer or binary values. Industrial MILP solvers such as Gurobi exploit the linear structure and the combinatorial structure of integer variables to produce effective exact and heuristic methods. However, many physically- and quantum-inspired methods operate natively on unconstrained, often nonlinear cost landscapes -- most notably the Quadratic Unconstrained Binary Optimization (QUBO) formulation, which is equivalent to a classically interacting spin model.

Converting a MILP into QUBO usually proceeds in two steps: (1) represent real-valued variables by binary arrays up to some chosen precision, and (2) turn constraints into soft quadratic penalty terms \cite{wang_quantum-inspired_2022}. That conversion is convenient for annealers and Ising-type machines, but it typically destroys the original problem structure that MILP solvers exploit.

This loss of structure matters because both MILP and QUBO formulations are NP-hard in general \cite{FBarahona_1982} and it is believed that no polynomial-time algorithm can guarantee optimality (or a bounded approximation ratio) in the worst case. Also, generic algorithms perform evenly across a class of problems (e.g.\ NP-hard problems), and only by exploiting structure, e.g.\ the one MILP solvers use, one can expect systematic gains \cite{NFL,ronnow2014defining}. In practice, then, a straightforward conversion of a carefully structured MILP into a penalized QUBO often yields instances on which classical solvers perform poorly -- not because the solvers are weak in general, but because the conversion removed exploitable structure and altered the problem’s landscape.

That observation motivates the search for cases where quantum or quantum-inspired devices may have an advantage. Several studies show that when a problem or a portion of it maps to a QUBO with features favourable to annealing (for example sparse interactions, locally hard substructures, or instances that admit advantageous minor embeddings) quantum algorithms can outperform a given classical algorithm on the QUBO formulation \cite{ajagekar_quantum_2019,ajagekar_quantum_2020,abbas_quantum_2023}.

Still, the literature is cautious: demonstrations to date document ``potential'' or ``limited quantum advantage'' rather than a general, hardware-agnostic breakthrough \cite{hoefler2023disentanglinghypepracticalityrealistically,ronnow2014defining}. The distinction is important and operational: potential quantum advantage is defined as a speedup relative to a specific classical algorithm or set of classical algorithms, and limited quantum advantage as a speedup relative to classical algorithms that implement the same algorithmic approach on classical hardware.

Practical benchmarking also exposes important caveats. Embedding overheads -- the cost of mapping problem graphs to the sparse connectivity of current QPUs -- are significant and can dominate runtimes, constraining effective problem size and eroding any raw speedup from the QPU itself \cite{ReinholdssonOdelius2023,naghmouchi_mixed_2024}. On the other hand, hybrid quantum-classical strategies -- where a classical decomposition exposes structure and a quantum component tackles the hard subproblems -- appear especially promising; they leverage classical pre/post-processing to preserve exploitable structure while using quantum resources where they are most likely to help \cite{franco_efficient_2023,VallejoBenitezCano2021}.

Taken together, these points argue for a cautious, structure-aware approach. Rather than expecting a blanket advantage, one should look for the structural components of an application that are classically challenging and naturally expressed as QUBO subproblems; those are the places where potential or limited quantum advantage may realistically appear \cite{ajagekar_quantum_2020,ronnow2014defining}. For domain problems such as food production -- where bottlenecks are combinatorial, high-dimensional, and time-sensitive -- hybrid quantum-classical pipelines can therefore represent a research-driven opportunity: they may accelerate decision-making and reduce decomposition costs in ways that pure classical solvers cannot, provided embedding and modeling overheads are managed appropriately.

\section{Implementation of a proof of concept}

% Instructions to be deleted before submission
\textit{
A proof of concept can vary in meaning depending on whether the proposed quantum computing solution is feasible on near-term hardware or requires fault-tolerant quantum computers anticipated in the longer term.
A proof of concept on near-term hardware might demonstrate that a simplified version of an algorithm can run on current noisy intermediate-scale quantum devices. 
A proof of concept on anticipated fault-tolerant quantum hardware might explore the scalability or theoretical performance of an algorithm.
}

\textit{
If not already clear from the previous section, please make clear what your proof of concept is about. 
If your proposal includes both a near-term implementation and an anticipated fault-tolerant implementation, please provide the relevant information in the below subsections 3.1-3.4 for both approaches.
}

\subsection{Description of the proof of concept}

% Instructions to be deleted before submission
\textit{
Please describe your proposed proof of concept to demonstrate the feasibility of the proposed quantum solution. This includes defining the regime of parameters to achieve quantum advantage and the problem size to be tackled (now and/or with future quantum computing devices). Please also specify whether access to real-world datasets has been secured and whether other known necessary initial conditions for the proof of concept are met. As always, include relevant references.
}


Our Proof of Concept consists of two main components: \textbf{scenarios} and \textbf{solvers}.

The solvers are structured as follows:
\begin{enumerate}
    \item Classical Solvers\begin{itemize}
        \item Linear problems are modeled via PuLP, which offers -- among others -- these two solvers:
        \begin{itemize}
            \item Gurobi \cite{gurobi2023}
            \item CPLEX \cite{cplex2023} (used as fallback)
        \end{itemize}
        \item Non Linear problems are modeled via Pyomo \cite{bynum2021pyomo}, which offers -- among other solvers:
        \begin{itemize}
            \item Gurobi \cite{gurobi2023}
            \item SCIP \cite{BolusaniEtal2024OO} (used as fallback)
        \end{itemize}
    \end{itemize}
    \item \begin{itemize}
        \item Quantum Solvers
        \begin{itemize}
            \item Dwave Leap Hybrid Solver \cite{DwaveLeapHybrid} which -- among others -- contains:
            \begin{itemize}
                \item LeapHybridCQMSampler for Constrained Quadratic Model (CQM) solver
                \item LeapHybridBQMSampler for Binary Quadratic Model (CQM) solver, which is used to solve the QUBO form of the problems
            \end{itemize}
        \end{itemize}
    \end{itemize}
\end{enumerate}

Scenarios on the other hand were created with the following procedure:

\paragraph{Crops:}

The data were provided by GAIN \cite{GAIN} in three different datasets:
\begin{itemize}
    \item \path{LCA results per kg & NVS.xlsx}
    \item \path{NVS_12Apr2024.xlsx}
    \item \path{PricePer100NVS_Indonesia_3Sept2024.xlsx}
\end{itemize}


The values were processed as follows to obtain normalized scores:  

\begin{itemize}
\item 
    \begin{equation}
       \text{Nut. Score}=\frac{\text{nutritional\_score}}{\max(\text{nutrient\_density})}
    \end{equation}
\item 
    \begin{equation}
       \text{Nut. Density}=\frac{\text{nutritional\_score}}{\max(\text{nutrient\_density})}
    \end{equation}
    \item 
    \begin{equation}
       \text{Env. Impact}=\left(\frac{\text{NVS}}{100}\times\frac{\text{Primary\_production\_mPt\_100NVS}}{\max(\text{Primary\_production\_mPt\_100NVS})}\right)
    \end{equation}
    \item 
    \begin{equation}
       \text{Sustainability}=\left(\frac{\text{NVS}}{100}\times\frac{\text{TOTAL\_categories\_mPt\_100NVS}}{\max(\text{TOTAL\_categories\_mPt\_100NVS})}\right)^{-1}
    \end{equation}
    \item
    \begin{equation}
       \text{Affordability}=\left(\frac{\text{NVS}}{100}\times\frac{\text{Cost to achieve NVS score of 100}}{\max(\text{Cost to achieve NVS score of 100})}\right)^{-1}
    \end{equation}
\end{itemize}

The precise content of each score deserves clarification, as the environmental and sustainability objectives are grounded in specific Life Cycle Assessment (LCA) impact categories from the GAIN database:
\begin{itemize}
    \item \textbf{Environmental Impact} is derived from \texttt{Primary\_production\_mPt\_100NVS}, which aggregates the environmental damage of the \textit{farming (primary production) stage only}, expressed in millipoints per 100 units of Nutritional Value Score (NVS). The underlying LCA categories at this stage include: climate change, acidification, particulate matter formation, eutrophication, land use, fossil resource use, and water use. By minimizing this score, the optimizer penalizes crops with high environmental damage during cultivation.
    \item \textbf{Sustainability} is derived from the inverse of \texttt{TOTAL\_categories\_mPt\_100NVS}, which sums all eight LCA damage categories (climate change, acidification, particulate matter, eutrophication, land use, fossil resource use, water use, and other) across the \textit{entire product lifecycle}---from primary production through processing, packaging, distribution, retail, consumer use, and waste treatment. Maximizing the inverse favors crops whose full supply chain has low total environmental damage per unit of nutritional value.
    \item \textbf{Affordability} is derived from the inverse of the cost required to achieve an NVS score of 100, sourced from Indonesian market price data. Maximizing this score favors crops that deliver high nutritional value at low cost to the consumer---i.e.\ the optimizer selects \textit{affordable} crops, not necessarily high-market-value crops.
    \item \textbf{Nutritional Value} and \textbf{Nutrient Density} are both derived from the NVS database and capture distinct aspects of nutritional quality per food item.
\end{itemize}

This distinction between Environmental Impact (production-stage damage) and Sustainability (full-lifecycle damage) allows the optimizer to separately penalize farming-phase externalities and total supply-chain externalities, providing finer-grained control over which environmental costs are weighted.

To create a comprehensive dataset, the same food was found in all three files (manually due to different naming conventions) and added to a separate dataset with the normalized scores. Then all rows containing empty values were discarded and the columns aggregated.

The aggregation is due to the fact that these datasets are not representative of a single geographic area, and don't represent a detailed enough dataset when considered singularly.

The data for the complete scenario is represented below:
\begin{table}[ht]
\centering
\scriptsize
\label{tab:food_data}
\begin{tabular}{llrrrrr}

\textbf{Food\_Name} & \textbf{Food Group}      & \textbf{Nut. Val.} & \textbf{Nut. Den.} & \textbf{Sust.} & \textbf{Env. Imp.} & \textbf{Afford.} \\

Mango          & Fruits                & 0.4468 & 0.2458 & 0.0757 & 0.0044 & 0.0261 \\
Papaya         & Fruits                & 0.4754 & 0.2748 & 0.1778 & 0.0172 & 0.0398 \\
Orange         & Fruits                & 0.4707 & 0.2537 & 0.1280 & 0.0083 & 0.0254 \\
Banana         & Fruits                & 0.4195 & 0.1963 & 0.1140 & 0.0088 & 0.0801 \\
Guava          & Fruits                & 0.5156 & 0.3102 & 0.1791 & 0.0120 & 0.0570 \\
Watermelon     & Fruits                & 0.3111 & 0.0706 & 0.0833 & 0.0089 & 0.0152 \\
Apple          & Fruits                & 0.3710 & 0.0884 & 0.0776 & 0.0045 & 0.0133 \\
Avocado        & Fruits                & 0.4674 & 0.2455 & 0.0511 & 0.0026 & 0.0357 \\
Durian         & Fruits                & 0.4516 & 0.2483 & 0.0275 & 0.0016 & 0.0203 \\
Corn           & Starchy staples       & 0.3908 & 0.1535 & 0.1214 & 0.0113 & 0.4179 \\
Potato         & Starchy staples       & 0.4782 & 0.3053 & 0.1246 & 0.0113 & 0.0934 \\
Tofu           & Pulses, nuts, and seeds & 0.5211 & 0.3471 & 0.1052 & 0.0188 & 0.1026 \\
Tempeh         & Pulses, nuts, and seeds & 0.5391 & 0.3946 & 0.1115 & 0.0201 & 0.2248 \\
Peanuts        & Pulses, nuts, and seeds & 0.4650 & 0.4268 & 0.0546 & 0.0031 & 0.2678 \\
Chickpeas      & Pulses, nuts, and seeds & 0.5153 & 0.3286 & 0.1404 & 0.0125 & 0.3980 \\
Pumpkin        & Vegetables            & 0.5889 & 0.4766 & 0.0579 & 0.0030 & 0.0338 \\
Spinach        & Vegetables            & 0.9032 & 0.9346 & 0.0859 & 0.0043 & 0.0362 \\
Tomatoes       & Vegetables            & 0.5816 & 0.4394 & 0.1039 & 0.0061 & 0.0387 \\
Long bean      & Vegetables            & 0.5616 & 0.4127 & 0.0821 & 0.0047 & 0.3634 \\
Cabbage        & Vegetables            & 0.6376 & 0.5007 & 0.0791 & 0.0043 & 0.0341 \\
Eggplant       & Vegetables            & 0.3967 & 0.1731 & 0.0597 & 0.0035 & 0.0217 \\
Cucumber       & Vegetables            & 0.4306 & 0.2272 & 0.1058 & 0.0084 & 0.0188 \\
Egg            & Animal-source foods   & 0.5837 & 0.4851 & 0.0343 & 0.0017 & 0.0217 \\
Beef           & Animal-source foods   & 0.5968 & 0.5424 & 0.0038 & 0.4468 & 0.0241 \\
Lamb           & Animal-source foods   & 0.5941 & 0.5332 & 0.0088 & 0.0005 & 0.0242 \\
Pork           & Animal-source foods   & 0.5840 & 0.5233 & 0.0165 & 0.0008 & 0.3743 \\
Chicken        & Animal-source foods   & 0.5533 & 0.4336 & 0.0249 & 0.0013 & 0.0572 \\

\end{tabular}
\caption{Food Data Overview}
\end{table}

\newpage
\paragraph{Farms:}

The farms were added to each scenario by sampling this distribution, obtained from \cite{LOWDER201616}.

\begin{table}[h!]
\centering

\begin{tabular}{lcc}

\textbf{Size class (ha)} & \textbf{Share of farms (\%)} & \textbf{Share of land (\%)} \\

$<1$        & $\sim 45$ & $\sim 10$ \\
$1$--$2$    & $\sim 20$ & $\sim 10$ \\
$2$--$5$    & $\sim 15$ & $\sim 20$ \\
$5$--$10$   & $\sim 8$  & $\sim 15$ \\
$10$--$20$  & $\sim 5$  & $\sim 20$ \\
$>20$       & $\sim 7$  & $\sim 25$ \\

\end{tabular}
\caption{Distribution of farm sizes and land shares in low-mid to low income countries}
\label{tab:farm_size_distribution}
\end{table}


\paragraph{Weights:}

The weights were chosen to be:



\begin{align*}
w_{\text{nutritional value}} &= 0.25 \\
w_{\text{nutrient density}} &= 0.2 \\
w_{\text{environmental impact}} &= 0.25 \\
w_{\text{affordability}} &= 0.15 \\
w_{\text{sustainability}} &= 0.15
\end{align*}

to maintain L-1 Normalization:

\begin{equation}
    \sum_i |w_i| =1
\end{equation}



\subsection{Benchmarking strategy}

\subsubsection{Scenario Adaptation}


Due to how the model is  formulated, the total area is a parameter known \textit{a priori}; this total area is going to get subdivided either into a number of farms with uneven area, or into a number of uniformly sized plots.

Additional information regarding the difference between the two choices can be found in the Methods \ref{sub:refinement}.

\subsubsection{Overview}
The benchmark will be comprehensive of two formulations:

\begin{enumerate}
    \item \textbf{Binary Formulation} (Even Grid): Used when land is divided into equal-sized plots
    \item \textbf{Continuous Formulation} (Uneven Distribution): Used when farms have varying sizes
\end{enumerate}

We aim to solve the optimization problem using three different methods:
\begin{itemize}
    \item PuLP with Gurobi solver (classical optimization)
    \item D-Wave Hybrid CQM Solver (quantum-classical hybrid)
    \item D-Wave Hybrid BQM Solver (quantum-enabled via CQM $\to$BQM conversion)
\end{itemize}



\subsubsection{Objective Functions}

\paragraph{Continuous Formulation Objective}

The objective function maximizes the weighted sum of agricultural value metrics, normalized by total available land:

$$\max \quad Z = \frac{1}{\sum_{f \in F} L_f} \sum_{f \in F} \sum_{c \in C} B_c \cdot A_{f,c}$$

where the composite value $B_c$ for crop $c$ is defined as:

$$B_c = w_{nv} \cdot v_{nv,c} + w_{nd} \cdot v_{nd,c} - w_{ei} \cdot v_{ei,c} + w_{af} \cdot v_{af,c} + w_{su} \cdot v_{su,c}$$

% \textbf{Note:} This is equivalent to \ref{eq:linear_obj} after appropriate renormalization.

\paragraph{Binary Formulation Objective}

For the binary formulation, the objective accounts for the fixed area $a_p$ of each plot:

$$\max \quad Z = \frac{1}{\sum_{p \in F} a_p} \sum_{p \in F} \sum_{c \in C} a_p \cdot B_c \cdot Y_{p,c}$$

where $B_c$ is defined identically as in the continuous formulation.

\textbf{Interpretation:} Each selected assignment contributes the plot's area multiplied by the crop's value density.

\subsubsection{Constraints}

\paragraph{Continuous Formulation Constraints}

\subparagraph{Land Availability Constraints}

Each farm cannot allocate more land than available:

$$\sum_{c \in \mathcal{C}} A_{f,c} \leq L_f \quad \forall f \in \mathcal{F}$$

\textbf{Label:} \texttt{Land\_Availability\_\{farm\}}

\subparagraph{Minimum Planting Area Constraints}

If a crop is selected on a farm, it must occupy at least the minimum required area:

$$A_{f,c} \geq A_{min,c} \cdot Y_{f,c} \quad \forall f \in \mathcal{F}, c \in C$$


\textbf{Logical Interpretation:}
\begin{itemize}
    \item If $Y_{f,c} = 1$: $A_{f,c} \geq A_{min,c}$ (enforces minimum area)
    \item If $Y_{f,c} = 0$: $A_{f,c} \geq 0$ (no planting, so area can be zero)
\end{itemize}

\textbf{Label:} \texttt{Min\_Area\_If\_Selected\_\{farm\}\_\{crop\}}

\subparagraph{Maximum Planting Area Constraints}

If a crop is not selected, its allocated area must be zero:

$$A_{f,c} \leq L_f \cdot Y_{f,c} \quad \forall f \in \mathcal{F}, c \in C$$


\textbf{Logical Interpretation:}
\begin{itemize}
    \item If $Y_{f,c} = 1$: $A_{f,c} \leq L_f$ (area can be up to farm size)
    \item If $Y_{f,c} = 0$: $A_{f,c} \leq 0$ (forces area to zero)
\end{itemize}

\textbf{Label:} \texttt{Max\_Area\_If\_Selected\_\{farm\}\_\{crop\}}

\subparagraph{Food Group Minimum Constraints}

At least a minimum number of \textit{unique} crops from specified food groups must be cultivated. Using the $U_c$ variables:

$$\sum_{c \in \mathcal{G}_g} U_c \geq N_{min,g} \quad \forall g \in \mathcal{G} \text{ where } N_{min,g} \text{ is defined}$$

\textbf{Label:} \texttt{Food\_Group\_Min\_\{group\}}

\subparagraph{Food Group Maximum Constraints}

A maximum number of \textit{unique} crops from specified food groups must not be exceeded:

$$\sum_{c \in \mathcal{G}_g} U_c \leq N_{max,g} \quad \forall g \in \mathcal{G} \text{ where } N_{max,g} \text{ is defined}$$

\textbf{Label:} \texttt{Food\_Group\_Max\_\{group\}}

\subparagraph{$U$-$Y$ Linking Constraints}

The $U_c$ variables must be properly linked to the $Y_{f,c}$ variables:

\begin{align}
Y_{f,c} &\leq U_c \quad \forall f \in \mathcal{F}, c \in \mathcal{C} \label{eq:u_link_upper}\\
U_c &\leq \sum_{f \in \mathcal{F}} Y_{f,c} \quad \forall c \in \mathcal{C} \label{eq:u_link_lower}
\end{align}

\textbf{Interpretation:}
\begin{itemize}
    \item Eq.~\eqref{eq:u_link_upper}: If any $Y_{f,c} = 1$, then $U_c$ must be 1
    \item Eq.~\eqref{eq:u_link_lower}: If no farm selects crop $c$ (all $Y_{f,c} = 0$), then $U_c$ must be 0
\end{itemize}

\textbf{Labels:} \texttt{U\_Link\_\{farm\}\_\{crop\}}, \texttt{U\_Bound\_\{crop\}}

\paragraph{Binary Formulation Constraints}

\subparagraph{Plot Assignment Constraints}

Each plot can be assigned to at most one crop (or remain idle):

$$\sum_{c \in \mathcal{C}} Y_{p,c} \leq 1 \quad \forall p \in \mathcal{F}$$



\textbf{Interpretation:}
\begin{itemize}
    \item $\sum_{c \in C} Y_{p,c} = 0$: Plot remains idle
    \item $\sum_{c \in C} Y_{p,c} = 1$: Plot is assigned to exactly one crop
\end{itemize}

\textbf{Label:} \texttt{Max\_Assignment\_\{plot\}}

\subparagraph{Minimum Plots Per Crop Constraints}

For crops with minimum planting area requirements, the constraint is converted to a minimum number of plots:

$$\sum_{p \in \mathcal{F}} Y_{p,c} \geq \left\lceil \frac{A_{min,c}}{a_p} \right\rceil \quad \forall p \in \mathcal{F} \text{ where } A_{min,c} > 0$$

where $a_p$ is the area of each plot (assumed equal in even grid).



\textbf{Interpretation:} If a crop $c$ requires minimum area $A_{min,c}$, it must be planted on at least $\lceil A_{min,c} / a_p \rceil$ plots.

\textbf{Label:} \texttt{Min\_Plots\_\{crop\}}

\subparagraph{Maximum Plots Per Crop Constraints}

For crops with maximum planting area limits, the constraint is converted to a maximum number of plots:

$$\sum_{p \in \mathcal{F}} Y_{p,c} \leq \left\lfloor \frac{A_{max,c}}{a_p} \right\rfloor \quad \forall c \in \mathcal{C} \text{ where } A_{max,c} \text{ is defined}$$



\textbf{Interpretation:} If a crop $c$ has maximum area $A_{max,c}$, it can be planted on at most $\lfloor A_{max,c} / a_p \rfloor$ plots.


\textbf{Label:} \texttt{Max\_Plots\_\{crop\}}




\subparagraph{Food Group Constraints}

The food group constraints use the $U_c$ variables to count \textit{unique} foods selected:

$$\sum_{c \in \mathcal{G}_g} U_c \geq N_{min,g} \quad \forall g \in \mathcal{G}$$
$$\sum_{c \in \mathcal{G}_g} U_c \leq N_{max,g} \quad \forall g \in \mathcal{G}$$

\subparagraph{$U$-$Y$ Linking Constraints}

The linking constraints ensure $U_c = 1$ if and only if crop $c$ is selected on at least one plot:

\begin{align}
Y_{p,c} &\leq U_c \quad \forall p \in \mathcal{F}, c \in \mathcal{C}\\
U_c &\leq \sum_{p \in \mathcal{F}} Y_{p,c} \quad \forall c \in \mathcal{C}
\end{align}

\textbf{Labels:} \texttt{MinFoodGroup\_Unique\_\{group\}}, \texttt{MaxFoodGroup\_Unique\_\{group\}}, \texttt{U\_Link\_\{plot\}\_\{crop\}}, \texttt{U\_Bound\_\{crop\}}


% \subsubsection{Plots}
% \clearpage


% \begin{figure}
%     \centering
%     \includegraphics[width=\linewidth]{Plots/Benchmark/performance_comparison.pdf}
%     %\caption{Caption}
%     %\label{fig:placeholder}
% \end{figure}


% \begin{figure}
%     \centering
%     \includegraphics[width=\linewidth]{Plots/Benchmark/solution_quality_comparison.pdf}
%     %\caption{Caption}
%     %\label{fig:placeholder}
% \end{figure}

% \clearpage    
% \begin{figure}
%     \centering
%     \includegraphics[width=\linewidth]{Plots/Benchmark/solution_composition_histograms.pdf}
%     %\caption{Caption}
%     %\label{fig:placeholder}
% \end{figure}


% \begin{figure}
%     \centering
%     \includegraphics[width=\linewidth]{Plots/Benchmark/solution_composition_pies.pdf}
%     %\caption{Caption}
%     %\label{fig:placeholder}
% \end{figure}

% }

\clearpage


\subsection{Resource estimation}

Resource estimation for quantum annealing applications differs fundamentally from gate-based quantum computing. Rather than circuit depth and gate counts, the relevant metrics center on problem embedding, qubit utilization, annealing time, and classical preprocessing overhead. This section provides comprehensive resource estimates for our crop allocation optimization across multiple problem scales and solution architectures.

\subsubsection{Problem Scaling and Variable Requirements}

The number of binary decision variables scales with the product of farms (or patches) and crops. For the binary crop allocation formulation, the total variable count is:
\begin{equation}
N_{\text{vars}} = |\mathcal{F}| \times |\mathcal{C}| + |\mathcal{C}|
\end{equation}
where the first term represents farm-crop assignment variables $Y_{f,c}$ and the second term accounts for global crop selection indicators $U_c$. For our benchmark configuration with 27 crops, this yields approximately $27|\mathcal{F}| + 27$ variables.

\begin{table}[H]
\centering
\caption{Problem scaling characteristics across benchmark configurations}
\label{tab:problem_scaling}
\begin{tabular}{@{}lccccc@{}}
\toprule
\textbf{Configuration} & \textbf{Farms} & \textbf{Crops} & \textbf{Variables} & \textbf{Constraints} & \textbf{Density} \\
\midrule
Validation & 10 & 27 & 297 & $\sim$120 & 0.15 \\
Small-scale & 50 & 27 & 1,377 & $\sim$350 & 0.08 \\
Medium-scale & 200 & 27 & 5,427 & $\sim$950 & 0.04 \\
Large-scale & 1,000 & 27 & 27,027 & $\sim$3,500 & 0.02 \\
\bottomrule
\end{tabular}
\end{table}

The constraint density decreases with problem size because the number of global constraints (food group diversity, crop area bounds) remains fixed while farm-level constraints scale linearly. This sparsity pattern is favorable for quantum annealing, as it limits the connectivity requirements in the embedded problem graph.

\subsubsection{Quantum Hardware Requirements}

\paragraph{D-Wave Advantage System Specifications}

Our target platform is the D-Wave Advantage quantum annealer, featuring 5,760 qubits in a Pegasus P16 topology with average connectivity of 15 neighbors per qubit. The Pegasus architecture supports native cliques of 15--20 fully connected qubits without chain overhead, enabling efficient embedding of moderately dense subproblems.

\begin{table}[H]
\centering
\caption{QPU resource utilization estimates}
\label{tab:qpu_resources}
\begin{tabular}{@{}lcccc@{}}
\toprule
\textbf{Problem Size} & \textbf{Logical Qubits} & \textbf{Physical Qubits} & \textbf{Chain Length} \\
\midrule
$\leq$50 vars & 50 & 80--150 & 1.2--1.8 \\
50--200 vars & 200 & 400--800 & 1.5--2.5  \\
200--500 vars & 500 & 1,200--2,500 & 2.0--4.0  \\
500--1,000 vars & 1,000 & 3,000--5,500 & 3.0--6.0  \\
\bottomrule
\end{tabular}
\end{table}

The chain length column indicates the average number of physical qubits required to represent each logical variable after minor embedding. Chain breaks occur when physical qubits within a chain disagree, directly impacting solution quality. Our empirical tests show chain break rates consistently below 2\% with auto-scaled chain strength, validating that the Pegasus topology accommodates our problem structure effectively.

% \paragraph{Annealing Time and Sample Requirements}

% Pure QPU annealing time scales favorably with problem size:
% \begin{equation}
% T_{\text{QPU}} \approx N_{\text{reads}} \times (t_{\text{anneal}} + t_{\text{readout}})
% \end{equation}
% where $t_{\text{anneal}} \in [20, 200]\,\mu s$ is the programmable annealing time and $t_{\text{readout}} \approx 100\,\mu s$ is the measurement overhead. For 100--500 reads per subproblem, total QPU access time remains under 100 ms even for large instances.

% \begin{table}[H]
% \centering
% \caption{Runtime estimates by solution architecture}
% \label{tab:runtime_estimates}
% \begin{tabular}{@{}lccc@{}}
% \toprule
% \textbf{Problem Size} & \textbf{Pure QPU} & \textbf{Hybrid Solver} & \textbf{Classical (Gurobi)} \\
% \midrule
% Small ($\leq$50 vars) & 2--100 ms & 10--60 s & $<$0.1 s \\
% Medium (50--200 vars) & 10--200 ms & 30--180 s & 0.1--1 s \\
% Large (200--500 vars) & 20--500 ms & 60--600 s & 1--10 s \\
% Very Large (500+ vars) & 50--1,000 ms & 300--3,600 s & 10--300 s \\
% \bottomrule
% \end{tabular}
% \end{table}

% The apparent paradox -- QPU times orders of magnitude faster than hybrid or classical solvers -- reflects that pure QPU time excludes embedding overhead and post-processing. Total wall-clock time for QPU-based solutions is dominated by classical preprocessing (embedding computation: 95--99\% of runtime) rather than quantum annealing itself \cite{reinholdsson2023embedding, mcgeoch2020dwave}.

\subsubsection{Empirical Simulated Annealing Benchmarks}

To ground our resource estimates in empirical data, we conducted simulated annealing (SA) experiments on both problem formulations using D-Wave's Ocean SDK (\texttt{SimulatedAnnealingSampler}) \cite{dwave_ocean} on a standard workstation (Intel i7-12700H, 16 GB RAM). Table~\ref{tab:sa_empirical} presents timing results across problem scales for the two formulations introduced in Section~\ref{sub:classical}.

\paragraph{BP (Binary Patch) Formulation}

The BP formulation uses pure binary variables $Y_{p,c} \in \{0,1\}$ and $U_c \in \{0,1\}$, making it directly compatible with QUBO/quantum annealing \cite{lucas2014ising, glover2018tutorial}. Total variables scale as 
\begin{equation}
N_{\text{vars}} = |\mathcal{F}| \times |\mathcal{C}| + |\mathcal{C}|
\end{equation}


\paragraph{MIP (Farm-level) Formulation with Discretization}

The original MIP formulation with continuous area variables $A_{f,c}$ cannot be directly solved by SA. To enable QUBO conversion, we discretize areas into $k$ binary levels \cite{chancellor2019domain}: $$A_{f,c} \approx \sum_{l=1}^{k} (L_f/k) \cdot a_{f,c,l} \quad\text{ where }a_{f,c,l} \in \{0,1\}$$ This significantly increases variable count to 

\begin{equation}
N_{\text{vars}} = |\mathcal{F}| \times |\mathcal{C}|\times k + |\mathcal{F}| \times |\mathcal{C}| + |\mathcal{C}|
\end{equation}

\begin{table}[H]
\centering
\caption{Empirical SA runtime measurements (100 reads, 1000 sweeps)}
\label{tab:sa_empirical}
\begin{adjustbox}{width=\textwidth}
\begin{tabular}{@{}lccccc@{}}
\toprule
\textbf{Formulation} & \textbf{Scale} & \textbf{Variables} & \textbf{Interactions} & \textbf{SA Time (s)} & \textbf{ms/var} \\
\midrule
\multirow{6}{*}{BP (Binary Patch)} 
 & 3 patches & 108 & 1,134 & 0.18 & 1.67 \\
 & 10 patches & 297 & 3,780 & 0.58 & 1.95 \\
 & 36 patches & 999 & 13,608 & 2.00 & 2.00 \\
 & 115 patches & 3,132 & 43,470 & 6.11 & 1.95 \\
 & 370 patches & 10,017 & 139,860 & 21.88 & 2.18 \\
 & 1,170 patches & 31,617 & 442,260 & 87.85 & 2.78 \\
\midrule
\multirow{4}{*}{MIP (4-level discretization)} 
 & 3 farms & 432 & 17,739 & 1.15 & 2.65 \\
 & 10 farms & 1,377 & 59,130 & 3.75 & 2.72 \\
 & 25 farms & 3,402 & 147,825 & 11.07 & 3.25 \\
 & 50 farms & 6,777 & 295,650 & 26.21 & 3.87 \\
\bottomrule
\end{tabular}
\end{adjustbox}
\end{table}

\paragraph{Scaling Analysis}

Fitting power-law models $T = \alpha \cdot n^\beta$ to the empirical data yields:
\begin{itemize}
    \item \textbf{BP formulation:} $T \approx 0.0012 \cdot n^{1.07}$ (nearly linear scaling)
    \item \textbf{MIP discretized:} $T \approx 0.0011 \cdot n^{1.14}$ (slightly super-linear due to denser interactions)
\end{itemize}

The BP formulation is more efficient per variable because it avoids the additional linking constraints required to enforce discretized area consistency. This demonstrates why the BP formulation is preferred for QUBO conversion: it produces smaller, sparser problems while preserving the essential optimization structure.

\paragraph{Comparison with Gurobi}

For context, Gurobi solves the \emph{original} MIP formulation (with continuous $A_{f,c}$ variables) in under 1 second for problems up to 1,000 farms, exploiting the tight LP relaxation \cite{gurobi2023}. However, when Gurobi solves the BP formulation as a pure binary MIP, performance degrades significantly -- the same problem that takes 0.1s with MIP may require 10--60s with BP due to the exponential worst-case complexity of branch-and-bound on pure binary problems without strong LP bounds \cite{wolsey1998integer, achterberg2007constraint}.

\subsubsection{Hybrid Solver Scaling from Literature}

Recent benchmarking studies provide empirical data on D-Wave hybrid solver performance for comparable QUBO/Max-Cut problems. Table~\ref{tab:hybrid_literature} summarizes key findings from Hrga et al.~\cite{hrga2024maxcut} and the KTH thesis by Reinholdsson \& Odelius~\cite{reinholdsson2023embedding}.

\begin{table}[H]
\centering
\caption{Hybrid solver runtime characteristics from literature}
\label{tab:hybrid_literature}
\begin{tabular}{@{}lccc@{}}
\toprule
\textbf{Problem Size} & \textbf{Hybrid Time (s)} & \textbf{SA-Quality Time (s)} & \textbf{Speedup} \\
\midrule
800 nodes & $\sim$3.0 & $\sim$672 & 224$\times$ \\
2,000 nodes & $\sim$5.2 & $\sim$670 & 129$\times$ \\
10,000 nodes & $\sim$5.2 & $\sim$680 & 131$\times$ \\
\bottomrule
\end{tabular}
\end{table}

Key observations:
\begin{itemize}
    \item The hybrid solver maintains \textbf{near-constant runtime} (3--5 seconds) across problem sizes from 800 to 10,000 variables, demonstrating excellent scaling properties \cite{hrga2024maxcut}.
    \item For specific problems below 251 variables, both hybrid and SA consistently achieve \textbf{global optima} \cite{hrga2024maxcut}.
    \item For larger problems (G-set Max-Cut benchmarks \cite{gset_benchmarks}), hybrid solver solutions are within 1--2\% of best-known values while being 100--200$\times$ faster than quality-focused SA \cite{hrga2024maxcut}.
    \item The Toshiba Simulated Bifurcation Machine (SBM) provides competitive quality at similar speeds to hybrid solvers \cite{goto2019combinatorial, toshiba_sbm}.
\end{itemize}

\paragraph{D-Wave Hybrid Solver Minimum Time Formula}

The LeapHybridBQMSampler uses piecewise-linear minimum runtime scaling \cite{dwave_hybrid_solver}:
\begin{equation}
T_{\text{min}}(n) = \begin{cases}
3\,\text{s} & n \leq 1024 \\
3 + 0.0023(n -- 1024)\,\text{s} & 1024 < n \leq 4096 \\
10 + 0.005(n -- 4096)\,\text{s} & 4096 < n \leq 10000 \\
\end{cases}
\end{equation}

For CQM problems with constraints, the formula incorporates additional terms \cite{dwave_cqm_solver}:
\begin{equation}
T_{\text{CQM}} = 5 + 0.000157 \cdot n_{\text{vars}} + 4.66 \times 10^{-6} \cdot n_{\text{biases}} + 6.44 \times 10^{-9} \cdot n_{\text{constraints}}
\end{equation}

\subsubsection{Comparative Runtime Projections}

Combining our empirical SA measurements with literature-derived hybrid scaling, we project comparative runtimes for our crop allocation problems:

\begin{table}[H]
\centering
\caption{Projected runtime comparison: SA vs Hybrid vs Gurobi}
\label{tab:runtime_comparison}
\begin{tabular}{@{}lcccc@{}}
\toprule
\textbf{Problem Instance} & \textbf{Variables} & \textbf{SA (ours)} & \textbf{Hybrid (proj.)} & \textbf{Gurobi} \\
\midrule
\multicolumn{5}{l}{\textit{BP Formulation (QUBO-native)}} \\
36 patches $\times$ 27 crops & 999 & 2.0 s & 3.0 s & 1--5 s \\
115 patches $\times$ 27 crops & 3,132 & 6.1 s & 3.0 s & 5--30 s \\
370 patches $\times$ 27 crops & 10,017 & 21.9 s & 5.2 s & 30--300 s \\
1,170 patches $\times$ 27 crops & 31,617 & 87.9 s & $\sim$8 s & timeout \\
\midrule
\multicolumn{5}{l}{\textit{MIP Formulation (original problem, Gurobi only)}} \\
100 farms $\times$ 27 crops & $\sim$2,800 & N/A & N/A & $<$0.5 s \\
500 farms $\times$ 27 crops & $\sim$14,000 & N/A & N/A & $<$1.0 s \\
1,000 farms $\times$ 27 crops & $\sim$28,000 & N/A & N/A & 1--5 s \\
\bottomrule
\end{tabular}
\end{table}

Note: The MIP formulation with continuous $A_{f,c}$ variables cannot be directly solved by SA or quantum annealing. Gurobi's performance on MIP is orders of magnitude better than on BP due to the tight LP relaxation.



\subsubsection{Energy Consumption Estimates (why not?)}

\paragraph{Quantum Annealing Energy Profile}

D-Wave systems operate at cryogenic temperatures ($\sim$15 mK), requiring continuous refrigeration \cite{johnson2011quantum}. The Advantage system consumes approximately 25 kW of power \cite{dwave_sustainability, fellous2022limitations}, dominated by the dilution refrigerator and classical control electronics rather than the quantum processor itself.

\begin{table}[H]
\centering
\caption{Energy consumption estimates by implementation scenario}
\label{tab:energy_estimates}
\begin{tabular}{@{}lccc@{}}
\toprule
\textbf{Scenario} & \textbf{Runtime} & \textbf{Energy (kWh)} & \textbf{CO$_2$ (kg)} \\
\midrule
PoC (100 experiments) & 5--10 hours & 125--250 & 50--100 \\
Production (daily) & 0.5--2 hours & 12.5--50 & 5--20 \\
Full-scale (weekly batch) & 10--20 hours & 250--500 & 100--200 \\
\bottomrule
\end{tabular}
\end{table}

Note: CO$_2$ estimates assume 0.4 kg CO$_2$/kWh (global average grid intensity). D-Wave's cloud infrastructure may use lower-carbon energy sources.

\paragraph{Comparison with Classical Computing}

For equivalent problem sizes, classical optimization on a high-performance workstation (500 W) requires:
\begin{itemize}
    \item Small problems: $<$1 second $\rightarrow$ negligible energy
    \item Large problems: 1--5 minutes $\rightarrow$ 0.004--0.04 kWh
\end{itemize}

The energy overhead of quantum annealing (due to cryogenic infrastructure) exceeds classical computing by 2--3 orders of magnitude for current problem sizes. This ratio improves for larger problems where QPU time remains nearly constant while classical time grows polynomially or exponentially.

\subsubsection{Hardware Provider}



D-Wave Leap provides the most mature quantum annealing infrastructure \cite{dwave_leap} with:
\begin{itemize}
    \item \textbf{Advantage system:} 5,760 qubits, Pegasus topology, ideal for QUBO/BQM problems \cite{boothby2020pegasus}
    \item \textbf{Hybrid solvers:} LeapHybridCQMSampler handles constrained problems up to 1 million variables \cite{dwave_hybrid_solver}
    \item \textbf{Access model:} Pay-per-use QPU time or monthly subscription plans
    \item \textbf{SDK:} Ocean toolkit with Python bindings for problem construction and embedding \cite{dwave_ocean}
\end{itemize}



% \subsubsection{Summary and Budget Recommendations}

% For a comprehensive proof-of-concept demonstration:

% \begin{itemize}
%     \item \textbf{Minimum viable:} 1 hour QPU time ($\sim$60--120 small experiments)
%     \item \textbf{Recommended:} 5 hours QPU time ($\sim$300 experiments across scales)
%     \item \textbf{Comprehensive benchmarking:} 20 hours QPU time (statistical analysis, parameter sweeps)
% \end{itemize}

% The hybrid solver approach (LeapHybridCQMSampler) requires no explicit QPU time budget but is billed separately. For constrained optimization problems like crop allocation, hybrid solvers provide the most robust path to solution quality while pure QPU experiments offer insight into annealing dynamics and decomposition strategies.



\subsection{Steps to achieve a proof of concept}

% Instructions to be deleted before submission
\textit{
Please describe the steps to deliver a successful proof of concept, from code writing to optimizing the code to tests on simulators and QPUs, including specifying what is conducted by whom in the team.
}

The proof of concept follows a three-phase development roadmap, progressing from initial implementation through scaling validation to optimization refinement. Each phase builds upon the previous, with explicit milestones and deliverables.

\subsubsection{Phase 1: Foundation and Classical Baseline}

\paragraph{Step 1.1: Data Pipeline Development}
Establish the data processing infrastructure to transform raw nutritional and environmental datasets into optimization-ready parameters:
\begin{itemize}
    \item Process GAIN nutritional databases to extract normalized scores for nutritional value, nutrient density, sustainability, environmental impact, and affordability
    \item Implement farm size sampling from the distribution in Table~\ref{tab:farm_size_distribution}
    \item Validate data consistency across all input sources
\end{itemize}

\paragraph{Step 1.2: Classical Solver Implementation}
Implement the MILP formulation using Gurobi as the classical baseline:
\begin{itemize}
    \item Primary solver: Gurobi via PuLP/Pyomo
    \item Establish baseline solution quality and runtime metrics across problem sizes
\end{itemize}

\paragraph{Step 1.3: Hybrid Solver Integration}
Convert the MILP to quantum-compatible formulations for D-Wave's hybrid solvers:
\begin{itemize}
    \item Constrained Quadratic Model (CQM): Native constraint handling via LeapHybridCQMSampler
    \item Binary Quadratic Model (BQM): QUBO transformation with penalty-based constraint encoding via LeapHybridBQMSampler
    \item Lagrange multiplier calibration study to identify optimal penalty strengths (see Section~\ref{sec:appx})
\end{itemize}

\subsubsection{Phase 2: Scaling Validation}

\paragraph{Step 2.1: Problem Size Progression}
Systematically increase problem complexity while maintaining solution quality:
\begin{itemize}
    \item Test instances: 5, 10, and 15 farms with 27 crops
    \item Monitor variable counts and constraint satisfaction as problems grow
\end{itemize}

\paragraph{Step 2.2: Simulator Testing}
Validate formulations using local simulation before committing hybrid solver resources:
\begin{itemize}
    \item Test CQM and BQM formulations using D-Wave's SimulatedAnnealingSampler
    \item Verify constraint satisfaction and solution feasibility across problem scales
\end{itemize}

\paragraph{Step 2.3: Comparative Benchmarking}
Execute rigorous statistical comparison between Gurobi and hybrid solvers:
\begin{itemize}
    \item Metrics: objective value, gap vs.\ Gurobi optimal (\%), wall-clock time, constraint violations
    \item Multiple runs per configuration to establish confidence intervals
\end{itemize}

\subsubsection{Phase 3: Analysis and Reporting}

\paragraph{Step 3.1: Large-Scale Validation}
Extend to operationally relevant problem sizes:
\begin{itemize}
    \item Target instances: 25, 50, and 100 farms
    \item Document solution quality and runtime across the full problem range
\end{itemize}

\paragraph{Step 3.2: Solution Quality Assessment}
Validate that hybrid solver solutions meet practical requirements:
\begin{itemize}
    \item Verify constraint satisfaction: land availability, minimum planting areas, food group diversity
    \item Compare objective quality against Gurobi baseline
    \item Characterize runtime behavior across problem scales
\end{itemize}

\subsubsection{Key Data Inputs and Requirements}

The proof of concept requires the following data categories:

\begin{enumerate}
\item \textbf{Population and Nutritional Needs}
\begin{itemize}
    \item Nutritional requirements: Calories, protein, dietary fats, essential micronutrients per population group
\end{itemize}

\item \textbf{Food System Parameters}
\begin{itemize}
    \item GAIN nLCA scores for nutrition-environment optimization
    \item Market prices for affordability scoring via Nutritional Value Score
\end{itemize}

\item \textbf{Production Constraints}
\begin{itemize}
    \item Total arable land availability and farm size distribution
    \item Agronomic compatibility constraints between crops
\end{itemize}
\end{enumerate}

\subsubsection{Deliverables and Success Criteria}

Upon completion of all phases, the proof of concept delivers:
\begin{itemize}
    \item Open-source solver framework supporting both classical (Gurobi) and hybrid (D-Wave) backends
    \item Benchmark dataset spanning 5--100 farm instances with documented solution quality
    \item Analysis report characterizing the performance of hybrid solvers relative to Gurobi across problem scales
\end{itemize}
\newpage
\section{Impact design}

% Instructions to be deleted before submission
\textit{
Anticipating a future deployment of the full-scale quantum solution in the real-world, please assess the impact of the proposed solution. An impact framework toolkit will soon be available to guide you through this assessment.
\begin{itemize}
\item At a local level, please assess the impact of your proposed solution on people, planet, and prosperity conducted across short-, medium-, and long-term timeframes. Detail the impacted areas and expected outcomes, documenting all assumptions and supporting references.
\item At a global level, how transposable is the use case to other contexts? What are other geographies or populations that would benefit from the proposed quantum solution (provide facts and references)? Would the conditions be conducive to a similar proof of concept (e.g. access to data)?
\item At a macro-level, please assess the interlinkages between SDGs to positively or negatively impact other SDGs (benefits vs risks), including mitigation measures if any?
\end{itemize}
}


\subsection{Local Impact Assessment}

The food production optimization solution targets multiple dimensions of sustainable development across different time horizons. As described in Section~1, both local and global populations are impacted, with particular vulnerability among low-income communities and regions with limited arable land. The solution addresses the fundamental optimization challenge of allocating limited farms and arable land to feed the maximum number of people in a nutritionally intelligent way, at the lowest cost to the population, while minimizing environmental harm. The optimizer pursues this through five explicit objectives---nutritional value, nutrient density, affordability, sustainability, and environmental impact---each grounded in quantitative data from LCA and nutritional databases (see Section~3.1 for definitions).

The anticipated impacts span nutritional security, environmental sustainability, and economic prosperity for agricultural communities.

\subsubsection{Short-Term Impacts (0-2 years)}

\textbf{SDG 2 (Zero Hunger):} The solution delivers optimized crop planning outputs that balance nutritional value, affordability, and environmental impact. Through quantum-classical hybrid optimization, the system produces planting assignments across diversified food groups that can be implemented immediately by farming communities. This addresses SDG Target 2.3 (doubling agricultural productivity of small-scale producers) and 2.4 (ensuring sustainable food production systems).

\textbf{SDG 1 (No Poverty):} By optimizing crop selection toward affordable, nutritionally dense foods, the system reduces food expenditure burdens on low-income households while maintaining dietary quality. Household income surveys can measure changes in the share of income spent on food compared to traditional planting patterns. This directly supports Target 1.2 (reducing poverty by at least half).

\textbf{SDG 9 (Industry, Innovation and Infrastructure):} The deployment of the digital mapping and quantum optimization platform provides stakeholders with accessible technology tools for agricultural planning, directly supporting SDG Target 9.a (facilitating sustainable infrastructure development in developing countries).

\textbf{SDG 10 (Reduced Inequalities):} The optimization tool is designed as open-source software whose outputs---crop allocation plans---can be communicated through existing agricultural extension services without requiring specialized hardware from farmers. Critically, the underlying farm model explicitly represents the fragmented land distribution typical of low-income countries: sub-1\,ha farms constitute approximately 45\% of holdings \cite{LOWDER201616}, and the optimizer handles these small parcels natively. This lowers the barrier for smallholder farmers to access data-driven decision support that was previously available only through costly private consultancy, addressing Target 10.2 (promoting social inclusion of all).

\subsubsection{Mid-Term Impacts (2-5 years)}

\textbf{SDG 2 (Zero Hunger):} Sustained implementation of optimized crop plans leads to measurable improvements in micronutrient intake. Biomarker analysis (blood/urine samples) can quantify changes in key nutrients (iron, vitamin A, zinc, folate) among participating households, addressing SDG Target 2.2 (ending all forms of malnutrition). Household Dietary Diversity Scores (HDDS) track the number of food groups consumed, demonstrating improved dietary variety.

\textbf{SDG 6 (Clean Water and Sanitation):} The sustainability objective explicitly includes \texttt{Water\_use\_mPt} as one of its eight LCA damage categories, meaning the optimizer penalizes crops with high water-use impact per unit of nutritional value. When consistently applied, this favors crop mixes with lower aggregate water footprints. Water meter readings and irrigation monitoring systems can quantify the resulting improvements in agricultural water consumption per season, addressing Target 6.4 (substantially increasing water-use efficiency).

\textbf{SDG 8 (Decent Work and Economic Growth):} As optimized crop diversification takes hold, additional employment opportunities emerge across sustainable agriculture---from planting through harvest and post-harvest processing. Employment surveys can document new labor positions generated by diversified farming operations, supporting Target 8.5 (full employment and decent work with equal pay).

\textbf{SDG 10 (Reduced Inequalities):} As smallholder farmers adopt the tool over successive planting seasons, the leveling effect of equal access to data-driven decision support may begin to narrow income gaps. However, we note that the optimizer does not directly target income inequality; rather, the mechanism is indirect---providing small-scale producers with crop plans of comparable quality to those available to larger, better-resourced operations. Monitoring changes in income distribution among participating farmers can track progress toward Target 10.1 (income growth of bottom 40\%).

\textbf{SDG 12 (Responsible Consumption and Production):} Farms adopting optimized production methods---guided by the environmental impact and sustainability objectives that penalize high-damage crops---increasingly shift toward environmentally sound practices. Practice surveys and certification records can document this transition, addressing Target 12.2 (sustainable management of natural resources).

\textbf{SDG 13 (Climate Action):} The environmental impact objective includes \texttt{Climate\_change\_mPt} at the production stage, while the sustainability objective captures climate change impact across the full lifecycle. When consistently applied, these objectives favor crop mixes with lower aggregate greenhouse gas emissions. GHG monitoring can quantify the resulting reductions in CO$_2$-equivalent emissions from agricultural activities, supporting climate mitigation efforts while maintaining food security.

\textbf{SDG 15 (Life on Land):} The sustainability objective includes \texttt{Land\_use\_mPt} as an LCA damage category, meaning the optimizer penalizes crops with high land-use impact per unit of nutritional value. This may contribute to reduced pressure on natural habitats by favoring land-efficient crops. Ecological surveys measuring species richness and abundance can document whether this translates to improved ecosystem health on participating farmland, addressing Target 15.5 (reducing degradation and halting biodiversity loss).

\textbf{SDG 17 (Partnerships for the Goals):} Successful adoption of the optimization solution requires multi-stakeholder partnerships between research institutions, government agencies, NGOs, and farmer cooperatives. Building these collaboration frameworks is a mid-term process; partnership agreements and MOUs documenting the resulting structures support Target 17.16 (enhanced global partnerships).

\subsubsection{Long-Term Impacts (5-10 years)}

\textbf{SDG 2 (Zero Hunger):} Long-term adoption leads to improved hunger rates at community and regional scales. Household dietary recall surveys and nutritional assessments show increased percentages of households meeting minimum dietary needs, addressing Target 2.1 (ending hunger and ensuring access to safe, nutritious, and sufficient food).

\textbf{SDG 4 (Quality Education) and SDG 8 (Decent Work):} Improved nutritional adequacy among communities translates to functional outcomes beyond hunger reduction: children show better school performance through improved attendance, concentration, and test scores (SDG~4, Target~4.1), while adults demonstrate improved work capacity and productivity (SDG~8, Target~8.5). These are downstream effects of sustained nutritional improvements rather than direct outputs of the optimizer.

\textbf{SDG 1 (No Poverty):} Poverty rates in farming communities decline measurably as sustained agricultural improvements compound over time. Economic surveys comparing household incomes to national poverty thresholds document the share of farming households moving above poverty lines, approaching Target 1.1 (eradicating extreme poverty).

\textbf{SDG 3 (Good Health and Well-being):} Reduced malnutrition and improved dietary diversity lower healthcare costs at the community level through decreased incidence of nutrition-related diseases. Average monthly household healthcare expenditure declines measurably, addressing Target 3.9 (reducing deaths and illnesses from pollution and contamination). This creates a positive feedback loop with agricultural productivity.

\textbf{SDG 8 (Decent Work and Economic Growth):} Agricultural sector contribution to regional GDP increases measurably as productivity gains scale across farming communities. Economic impact assessments document the sector's enhanced role in sustainable growth, supporting Target 8.1 (sustaining per capita economic growth).

\textbf{SDG 13 (Climate Action):} Cumulative emissions reductions from optimized agricultural systems contribute to global warming mitigation efforts, with measurable impacts on regional climate resilience. GHG monitoring tools, soil tests, and energy-use records document greenhouse gas reduction per farming season, supporting Target 13.2 (integrating climate measures into policies and planning).

\textbf{SDG 15 (Life on Land):} Reduced land-use intensity, as incentivized by the sustainability objective's inclusion of \texttt{Land\_use\_mPt}, may contribute to lower pressure on soils over time. However, direct soil quality improvement (organic matter, nutrient content, pH) would require complementary agronomic practices such as cover cropping and reduced tillage that are not explicitly modeled by the optimizer. Soil testing at baseline and regular intervals can track whether optimized crop selection correlates with improved soil indicators, addressing Target 15.3 (combating desertification and restoring degraded land).

\subsection{SDG Interlinkages: Synergies, Trade-offs, and Mitigation Strategies}

The quantum optimization approach creates a complex network of SDG interlinkages that must be carefully managed to maximize synergies and minimize trade-offs \cite{Fanzo2022,UNFSS_Planet_2023,Sylvester2024}. Understanding these relationships enables responsible innovation that advances holistic sustainable development rather than optimizing narrow metrics at the expense of broader welfare.

\subsubsection{Primary Synergies}

Several mutually reinforcing pathways amplify the solution's positive impacts across multiple SDGs simultaneously.

The \textbf{nutrition-health-productivity pathway} (SDG 2 $\rightarrow$ SDG 3 $\rightarrow$ SDG 8) creates a virtuous cycle. Optimizing for nutritional density directly improves health outcomes, which in turn enhances workforce productivity and agricultural capacity. This reinforcing loop compounds benefits over time as healthier populations demonstrate greater economic participation and better educational attainment.

The \textbf{poverty-food-economics triangle} interconnects poverty reduction (SDG 1), food security (SDG 2), and economic growth (SDG 8) in a mutually supporting relationship. Agricultural optimization addressing all three simultaneously achieves greater total impact than addressing any single goal in isolation. By maximizing affordability (low cost per unit of nutritional value), the optimizer reduces food expenditure burdens on consumers; combined with dietary diversification, this improves nutritional security while freeing household income for other needs, contributing to local economic activity.

The \textbf{technology-innovation-partnership chain} demonstrates how quantum computing application (SDG 9) strengthens international partnerships (SDG 17) by facilitating cutting-edge technology transfer to developing contexts while building local technical capacity. This pathway enables knowledge sharing, reduces inequalities through accessible technology (SDG 10), and drives economic growth (SDG 8) through innovation adoption.

The \textbf{climate-ecosystem nexus} shows how the environmental optimization objectives simultaneously address multiple environmental SDGs through their underlying LCA damage categories. Specifically: \texttt{Climate\_change\_mPt} directly relates to climate action (SDG~13); \texttt{Land\_use\_mPt} relates to terrestrial ecosystem pressure (SDG~15); \texttt{Water\_use\_mPt} relates to water management (SDG~6); and \texttt{Eutrophication\_mPt} captures nutrient runoff that can degrade aquatic systems (SDG~14). By penalizing crops with high scores in these categories, the optimizer indirectly favors crop mixes with lower aggregate environmental footprints across these dimensions. However, biodiversity protection (SDG~15) is only indirectly supported through reduced land-use pressure; the optimizer does not explicitly model species richness or habitat preservation.

\subsubsection{Critical Trade-offs and Mitigation Strategies}

While the solution generates numerous synergies, several fundamental tensions require active management to prevent negative impacts.

\textbf{Productivity-Environment Tension (SDG 2 vs SDGs 6, 13, 14, 15):}

The fundamental trade-off between maximizing food production and minimizing environmental impact represents the core challenge of sustainable agriculture. Without explicit constraints, optimization might select high-yield crops that deplete water resources, increase GHG emissions, harm terrestrial ecosystems, or generate agricultural runoff degrading aquatic environments.

\textit{Mitigation:} The optimization framework employs explicit multi-objective formulations that balance nutritional adequacy and affordability with environmental sustainability. Environmental impact scores from life-cycle assessments constrain solutions to minimize land use, water consumption, and carbon footprint while maintaining nutritional adequacy. Pareto analysis reveals optimal balance points where multiple objectives coexist rather than compete. Continuous monitoring ensures environmental constraints remain effective as agricultural practices evolve.

\textbf{Technology Access and Digital Divide (SDG 9 vs SDG 10):}

While quantum optimization provides powerful capabilities, it risks exacerbating inequalities if access is limited to well-resourced actors. Technology adoption may inadvertently widen the gap between large commercial operations and smallholder farmers, undermining poverty reduction (SDG 1) and equality (SDG 10) goals even while advancing innovation (SDG 9).

\textit{Mitigation:} Access equity requires deliberate design choices including: open-source software releases ensuring zero-cost availability; user-friendly interfaces requiring minimal technical expertise; partnerships with agricultural extension services providing last-mile technology dissemination to smallholder farmers; community-based facilitator models bridging the digital divide; differentiated support mechanisms for smaller farms; cooperative models enabling resource pooling; and progressive pricing structures favoring disadvantaged farmers. Capacity-building programs ensure technology access translates to effective utilization.

\textbf{Climate Policy Costs (SDG 13 vs SDGs 1, 2, 8):}

Environmental optimization may increase short-term costs for implementing sustainable practices, potentially harming poor households and slowing immediate economic growth. The transition to climate-smart agriculture requires upfront investments that may be prohibitive for resource-constrained farmers, creating tension between climate action and poverty reduction.

\textit{Mitigation:} Balancing climate goals with poverty reduction and food security requires complementary policy interventions: phased implementation allowing gradual adaptation; transition subsidies protecting vulnerable households during the shift to sustainable practices; technical assistance reducing implementation costs through knowledge transfer; market development for sustainable products commanding premium prices; payments for ecosystem services compensating farmers for environmental benefits they generate; social protection systems providing safety nets during transitions; climate-resilient livelihood diversification reducing farmer vulnerability; and staged optimization ensuring affordability remains a weighted objective across all time horizons. Just transition frameworks ensure equitable distribution of costs and benefits.

\textbf{Conservation-Production Balance (SDG 15 vs SDG 2):}

Protecting biodiversity and preventing land degradation may constrain agricultural expansion, creating tension between environmental protection and food security. In land-scarce contexts, preserving natural areas for ecosystem health can limit the area available for food production.

\textit{Mitigation:} Resolution requires landscape-level approaches including: sustainable intensification on existing farmland rather than extensification into natural areas; agroforestry and agroecological systems integrating conservation with production; landscape-level planning designating distinct conservation and production zones; restoration of degraded lands for agricultural use; payment schemes compensating farmers for conservation activities; and demonstrating that soil health improvements and ecosystem services support long-term agricultural productivity, aligning conservation with food security objectives.

\textbf{Healthcare Infrastructure Expansion (SDG 3 vs SDGs 12, 13):}

\textit{Note: While improved nutrition may increase demand for healthcare services in the long term as populations grow, this effect is several causal steps removed from the optimizer's direct outputs and is not considered a primary trade-off of the solution.}

\textbf{Economic Growth Without Safeguards (SDG 8 vs SDG 13):}

Agricultural economic growth, if pursued without environmental constraints, may increase emissions and resource consumption. The drive for productivity and income gains could undermine climate goals if environmental impacts are not explicitly managed.

\textit{Mitigation:} The optimization explicitly constrains environmental impacts while maximizing economic value, ensuring growth aligns with sustainability. Long-term strategies include mainstreaming environmental safeguards in agricultural policy, implementing carbon pricing or offset mechanisms, and demonstrating that sustainable intensification supports both growth and climate goals simultaneously.

\textbf{Consumption Patterns and Economic Growth (SDG 12 vs SDG 8):}

The sustainability objective includes \texttt{Waste\_treatment\_mPt} as part of the full-lifecycle environmental assessment, meaning the optimizer indirectly penalizes crops whose post-consumer waste processing has high environmental costs. However, the optimizer does not directly model consumer behavior, food waste rates, or consumption patterns. Claims about responsible consumption are therefore limited to the production side: the optimizer favors crops with lower full-lifecycle environmental damage, which includes the waste treatment stage. Promoting responsible consumption at the consumer level would require complementary interventions beyond the scope of this optimization.

\textit{Mitigation:} Market development for sustainable products commanding premium prices transforms environmental responsibility into economic opportunity. Certification programs add value to sustainably produced crops, while consumer education creates demand. Long-term economic benefits of sustainable practices---including reduced input costs and climate resilience---outweigh short-term transition costs.

\textbf{Unequal Partnership Structures (SDG 17 vs SDG 10):}

Multi-stakeholder partnerships, while essential for technology deployment, may reinforce existing power imbalances if not carefully structured. Partnerships dominated by well-resourced institutions could marginalize smallholder farmer voices. Furthermore, quantum computing hardware and cloud services are predominantly provided by technology companies in North America and Europe, which could create technological dependencies for developing countries relying on these platforms for agricultural decision support.

\textit{Mitigation:} Participatory governance structures with farmer representatives in decision-making bodies ensure equitable influence. Transparent benefit-sharing agreements document how gains are distributed. Capacity building programs equalize negotiating power between partners. Regular stakeholder consultations verify that partnership benefits reach intended beneficiaries. To mitigate technological dependency, the optimization framework is designed to also run on classical solvers (Gurobi, SCIP) when quantum access is unavailable, ensuring the solution remains functional independently of quantum cloud providers.

\subsubsection{Monitoring and Adaptive Management}

Responsible innovation requires ongoing assessment of both intended and unintended consequences across the SDG network. The impact framework includes: baseline and endline surveys for all indicator measurements, participatory monitoring engaging farming communities in data collection and interpretation, regular SDG interlinkage assessments using network analysis methods \cite{Fanzo2022} to identify emerging trade-offs or unexpected synergies, and adaptive optimization that adjusts objective weights based on observed outcomes and stakeholder feedback. This adaptive approach ensures the solution remains aligned with holistic sustainable development goals rather than inadvertently creating negative externalities while optimizing for specific metrics.

\subsection{Global Transposability}

The quantum-enabled crop optimization approach is highly transposable to diverse geographical contexts, particularly in regions facing the triple burden of malnutrition, environmental degradation, and climate vulnerability.

\subsubsection{Target Geographies}

\textbf{Sub-Saharan Africa:} With 45\% of the population in rural areas and less than 10\% of health professionals serving these communities \cite{Fanzo2022}, optimized agricultural planning can address both nutritional deficiencies and economic development. Small-scale farms ($<$1 hectare) constitute approximately 45\% of farms but control only 10\% of agricultural land \cite{LOWDER201616}, making optimization critical for efficient resource utilization.

\textbf{South and Southeast Asia:} High population density, diverse crop systems, and significant climate risks make these regions ideal candidates. The methodology has been validated using Indonesian food data, demonstrating cultural adaptability through locally acceptable food baskets.

\textbf{Latin America:} Regions with high biodiversity but also significant deforestation pressures (80\% driven by food systems \cite{Fanzo2022}) would benefit from optimization that explicitly balances productivity with ecosystem conservation.

\textbf{Small Island Developing States (SIDS):} Limited arable land and climate vulnerability necessitate highly efficient crop planning. The optimization approach can maximize nutritional output per hectare while minimizing import dependence.

\subsubsection{Data Requirements and Deployment Feasibility}

The proof of concept requires three data categories, all of which are increasingly accessible in diverse global contexts. First, nutritional composition and environmental impact scores are available from GAIN databases and life-cycle assessments, providing standardized inputs across geographies. Second, local food preferences and prices can be collected through household surveys and market data, ensuring cultural relevance and economic realism. Third, farm characteristics and constraints are accessible through agricultural census data or participatory mapping exercises.

These data types are becoming widely available through global initiatives like the Food Systems Dashboard, making deployment feasible in most contexts with appropriate partnerships. The modular architecture allows regions with partial data availability to implement simplified versions while building toward more comprehensive optimization as data infrastructure improves.

\subsection{Implementation Considerations}

\subsubsection{Key Risks}

Several risks could affect impact achievement across different time horizons and must be proactively managed. Income inequality between farms may persist or widen if optimization benefits apply primarily to well-resourced farmers rather than smallholders, affecting both short-term and mid-term outcomes. Lack of market access for new crops poses a significant barrier—even optimally selected crops cannot improve incomes if buyers are unavailable, particularly impacting short-term and mid-term economic gains. Regulatory barriers to new seed varieties or agricultural inputs may slow implementation at the short-term stage.

Climate-related risks span multiple timeframes. Crops may not be resilient enough for accelerating climate change impacts, threatening both short-term and mid-term productivity gains. More acutely, severe climate shocks (droughts, floods, extreme temperatures) can dramatically reduce yields across all timeframes—short-term, mid-term, and long-term—potentially undermining the entire impact logic chain. Finally, poor quality or missing baseline data compromises the ability to measure impacts accurately in the short term, weakening evidence generation and adaptive management.

\subsubsection{Critical Assumptions}

Several foundational assumptions underpin the impact pathway and require validation during implementation. First, local input supply chains must be able to scale to support new crop varieties—seed suppliers, fertilizer distributors, and agricultural extension services need sufficient capacity to meet changing demands. Second, buyers will purchase new nutrient-dense crops, creating market incentives for farmers to adopt optimized crop plans. Third, national regulations must permit trial of new varieties, enabling innovation while maintaining biosafety and quality standards. Fourth, farmers will accept and implement optimized crop plans if they prove profitable, assuming economic rationality drives adoption decisions. Fifth, sufficient soil and climate data are available for baseline surveys, enabling accurate optimization and impact measurement.

\subsubsection{Preparatory Activities}

Three essential activities lay the groundwork for successful implementation. Conducting baseline soil, climate, and household surveys establishes the data foundation for optimization algorithms and impact evaluation, quantifying initial conditions across agronomic, environmental, and socioeconomic dimensions. Linking producer groups with buyers and processors creates the market channels necessary for farmers to sell optimized crop outputs, addressing the market access risk proactively. Engaging regulators and documenting compliance needs ensures regulatory pathways are clear before implementation, reducing delays and enabling adaptive policy development that supports innovation while protecting public interests.

\subsection{Stakeholder Engagement}

Successful implementation requires coordinated engagement with diverse stakeholders, each playing distinct roles in the impact pathway. The Ministry of Agriculture serves as both regulator and policy architect, with primary interests in ensuring food security and regulatory compliance. As a direct stakeholder, engagement occurs through policy dialogue mechanisms, establishing regulatory frameworks that enable innovation while protecting public interests.

Farmers constitute the primary producers and direct beneficiaries of the optimization system. Their central interest—increasing yields and income while managing risk—drives adoption decisions. Direct engagement through field training programs builds capacity for using optimization tools and implementing recommended crop plans. Complementing farmer engagement, buyers and offtakers (including agricultural cooperatives, food processors, and wholesalers) require reliable supplies of nutritious crops. While indirect stakeholders, buyers are critical to closing the value chain. Market linkage activities connect optimized production with demand, ensuring farmers can sell diversified outputs profitably.

Retailers sell nutrient-dense crops to end consumers, with interests in maintaining reliable supply of diverse, high-quality produce. As direct stakeholders, retailers are engaged through market linkage strategies that align optimized crop production with retail demand patterns and quality standards. The local community represents the ultimate beneficiaries—end users of improved food systems seeking better nutrition, stable food access, and lower costs. Direct community engagement through outreach programs ensures the optimization system responds to local food preferences, affordability constraints, and cultural contexts, making nutrition improvements socially sustainable.






\section{Methods}
\label{sec:appx}




% Standalone subsection: derivation of the project model as a subset of Esteso et al. (2022)
% To be \input{} or pasted into the Methods section of the report.

\subsection{Derivation from the reference model}
\label{sub:model_derivation}

Our formulation is derived from the multi-objective sustainable crop planning model of Esteso et al.~\cite{esteso_sustainable_2023}, which we simplify along several dimensions to obtain a problem instance suitable for QUBO reformulation and quantum annealing benchmarks. We summarize below the full reference model and detail each reduction.

\subsubsection{The reference model}

Esteso et al.\ propose a multi-echelon supply-chain (SC) optimization model for sustainable crop production.
The model spans a multi-period planning horizon (indexed by~$t$) and covers three echelons---farms, a shared drying facility, and processing plants---connected by transport logistics.

\paragraph{Indices and sets.}
The reference model is defined over the following index sets: crops~$c$, farms~$f$, processing plants~$l$, planting periods~$p$, harvest periods~$h$, harvest seasons~$s$, and a growth-time index $z$ (the number of periods between planting or consecutive harvests).

\paragraph{Decision variables.}
The full model contains thirteen families of decision variables:
\begin{enumerate}
    \item $AP^{p}_{cf}$~(continuous): area planted with crop~$c$ on farm~$f$ in planting period~$p$;
    \item $AH^{shz}_{cf}$: area harvested, indexed by season~$s$, period~$h$, and growth time~$z$;
    \item $QH^{shz}_{cf}$, $QD^{shz}_{cf}$, $QCD^{shz}_{cf}$: quantities harvested, dried on-farm, and dried at the central facility;
    \item $ID^{shzt}_{cf}$, $QS^{shzt}_{cf}$, $IS^{shzt}_{cf}$: inventories of dried and sorted crops, and quantities sorted;
    \item $QT^{shzt}_{cfl}$: quantity transported from farm~$f$ to processing plant~$l$;
    \item $YT^{t}_{cfl}$~(binary): indicates whether a transport shipment occurs;
    \item $I^{t}_{cl}$: inventory at processing plants;
    \item $Pr_f$, $U_f$~(continuous): farm profit and unfairness measure.
\end{enumerate}

\paragraph{Objectives.}
Three objectives are optimized simultaneously via the $\varepsilon$-constraint method:
\begin{align}
    \max\; Z_A &= \sum_{c,f,l,s,h,z,t} am^{sz}_{c}\cdot QT^{shzt}_{cfl}
        & \text{(active-molecule concentration),} \label{eq:esteso_ZA}\\[4pt]
    \min\; Z_C &= \sum (\text{planting} + \text{harvesting} + \text{drying} \notag\\
               &\quad + \text{sorting} + \text{inventory} + \text{transport costs})
        & \text{(SC costs),} \label{eq:esteso_ZC}\\[4pt]
    \min\; Z_U &= \sum_f U_f
        & \text{(farmer unfairness).} \label{eq:esteso_ZU}
\end{align}
$Z_A$ maximizes the total concentration of active molecules in deliveries to processors, where $am^{sz}_c$ depends on the growth time~$z$. $Z_C$ minimizes the aggregate cost across five SC stages (planting, harvesting, drying, sorting, inventory, and transport). $Z_U$ minimizes economic unfairness among farmers, measured as the deviation of each farm's profit-per-hectare from the average.

\paragraph{Constraints.}
Sixteen constraint families enforce: (i)~farm area limits, (ii)~harvest-to-planting area balances across seasons, (iii)~yield equations linking harvest quantity to area and growth time, (iv)~on-farm and central drying capacity limits, (v)~inventory balance equations for dried and sorted products, (vi)~sorting and storage capacity limits, (vii)~minimum transport lot sizes with binary linking, (viii)~processor inventory balance and demand fulfillment, and (ix)~the linearization of the unfairness measure.

In the case study of the reference paper (10~farms, 2~crops, 104-week horizon), the resulting MILP instance contains 413\,844 decision variables and 292\,717 constraints.

\subsubsection{Reductions applied}
\label{sub:reductions}

We retain the general structure of centralized multi-farm crop allocation from the reference model, but apply the following simplifications to isolate the combinatorial core of the problem.

\paragraph{1. Elimination of the temporal dimension.}
The multi-period planning horizon is collapsed to a single static allocation period. All time indices ($p$, $h$, $s$, $t$) and the growth-time variable~$z$ are removed. As a consequence, the harvest-area balances (Eqs.~6--7 of the reference), the yield equations (Eq.~8), and all inventory balance constraints (Eqs.~12--13, 15, 19) are no longer applicable.

\paragraph{2. Elimination of post-harvest supply-chain stages.}
We restrict the model to the farm echelon. The drying (Eqs.~9--11), sorting (Eq.~14), storage (Eq.~16), and transport (Eqs.~17--18) stages are removed together with their associated variables ($QD$, $QCD$, $QS$, $QT$, $YT$, $ID$, $IS$, $I$) and capacity parameters. This eliminates the entire downstream SC structure and all inter-echelon flow constraints.

\paragraph{3. Replacement of the objective functions.}
The three original objectives---active-molecule concentration~\eqref{eq:esteso_ZA}, supply-chain costs~\eqref{eq:esteso_ZC}, and farmer unfairness~\eqref{eq:esteso_ZU}---are replaced by four objectives relevant to food security and dietary adequacy:
\begin{enumerate}
    \item \emph{maximize nutritional value} (score $v_{nv,c}$),
    \item \emph{maximize nutrient density} (score $v_{nd,c}$),
    \item \emph{maximize affordability} (score $v_{af,c}$),
    \item \emph{minimize environmental impact} (score $v_{ei,c}$).
\end{enumerate}
These per-crop scores are derived from external nutritional and environmental databases (see Section~\ref{...}) and are \emph{static}, in contrast to the growth-time-dependent molecule concentration of the reference model. In addition, a sustainability score~$v_{su,c}$ is included to capture the aggregate environmental footprint.

The $\varepsilon$-constraint resolution of the reference is replaced by a weighted-sum scalarization: the four (plus sustainability) objectives are combined into a single composite benefit
\begin{equation}
    B_c = w_{nv}\,v_{nv,c} + w_{nd}\,v_{nd,c} - w_{ei}\,v_{ei,c} + w_{af}\,v_{af,c} + w_{su}\,v_{su,c}\,,
    \label{eq:composite_benefit}
\end{equation}
yielding the scalar objective presented in Eq.~\eqref{eq:generalized objective}.

\paragraph{4. Simplification of the decision variables.}
Of the thirteen variable families in the reference model, only the planted-area variable $AP^{p}_{cf}$ is retained, in simplified form as $A_{f,c} \in [0, L_f]$ (continuous, without the temporal index). Two new binary variable families are introduced to support the MILP structure:
\begin{itemize}
    \item $Y_{f,c} \in \{0,1\}$: indicates whether food~$c$ is planted on farm~$f$;
    \item $U_c \in \{0,1\}$: indicates whether food~$c$ is planted on \emph{at least one} farm (used for counting distinct foods in food-group diversity constraints).
\end{itemize}

\paragraph{5. Retained and added constraints.}
From the sixteen constraint families in the reference model, only the farm area limit (Eq.~5) is retained:
\begin{equation}
    \sum_{c \in \mathcal{C}} A_{f,c} \leq L_f \quad \forall\, f \in \mathcal{F}\,.
    \label{eq:land_avail}
\end{equation}
Five new constraint families---absent from the reference model---are introduced to encode the agronomic and dietary requirements of the simplified problem:
\begin{itemize}
    \item \emph{Minimum planting area}: $A_{f,c} \geq A_{\min,c}\cdot Y_{f,c}$;
    \item \emph{Selection linking}: $A_{f,c} \leq L_f \cdot Y_{f,c}$;
    \item \emph{Food-group diversity} (min/max): $N_{\min,g} \leq \sum_{c \in \mathcal{G}_g} U_c \leq N_{\max,g}$;
    \item \emph{$U$--$Y$ linking}: $Y_{f,c} \leq U_c$ and $U_c \leq \sum_{f} Y_{f,c}$.
\end{itemize}

\paragraph{6. Change of application domain.}
The reference model targets biennial medicinal plants (lemon balm and sage) grown for active-molecule extraction. Our model addresses 27~food crops spanning five food groups (fruits, vegetables, starchy staples, pulses/nuts/seeds, and animal-source foods), selected for nutritional adequacy and dietary diversity in a development context.

\subsubsection{Summary of the subset}
Table~\ref{tab:model_reduction} summarizes the relationship between the reference model and our reduced formulation.

\begin{table}[ht]
\centering
\small
\begin{tabular}{lcc}
\toprule
\textbf{Model element} & \textbf{Esteso et al.} & \textbf{This work} \\
\midrule
Objectives             & 3  & 4 (+1 sust.) \\
Decision var.\ families & 13 & 3 \\
Constraint families    & 16 & 6 \\
Time periods           & 104 & 1 (static) \\
SC echelons            & 3  & 1 (farm only) \\
Crops                  & 2  & 27 \\
Resolution method      & $\varepsilon$-constraint + TOPSIS & weighted sum \\
\bottomrule
\end{tabular}
\caption{Comparison between the full reference model of Esteso et al.~\cite{esteso_sustainable_2023} and the reduced formulation used in this work. The reduction eliminates the temporal dimension and all post-harvest SC stages, retains the farm area constraint, replaces the original objectives with nutritional and environmental scores, and adds food-group diversity constraints.}
\label{tab:model_reduction}
\end{table}

The rationale for these simplifications is twofold. First, removing the temporal dynamics and the downstream supply chain isolates the \emph{combinatorial crop-allocation core} of the problem---the assignment of crops to farms subject to area and diversity constraints---which is the component most naturally amenable to QUBO reformulation and quantum annealing (see Section~\ref{PoC}). Second, replacing the domain-specific medicinal-plant objectives with broadly applicable nutritional and environmental scores allows the model to address the food-security context motivating this work, while preserving the multi-objective structure of the original formulation.


\subsection{Lagrange Multiplier Calibration Study}

\subsubsection{Motivation and Objective}

When converting a Constrained Quadratic Model (CQM) to a Binary Quadratic Model (BQM) via the penalty method, constraint violations are penalized using Lagrange multipliers $\lambda$. The choice of $\lambda$ involves a critical trade-off:

\begin{itemize}
    \item \textbf{Too small ($\lambda \to 0$):} Penalties are weak, leading to constraint violations in the final solution
    \item \textbf{Too large ($\lambda \to \infty$):} Penalties dominate the objective, suppressing optimization and potentially causing numerical instability
\end{itemize}

The goal of this study is to identify the \textbf{minimum Lagrange multiplier} that achieves zero constraint violations while preserving objective optimization quality.

\subsubsection{Experimental Design}

\textbf{Test Instance:}
\begin{itemize}
    \item Scenario: \texttt{full\_family} with 10 patches (small-scale for rapid testing)
    \item Patch generation: \texttt{generate\_farms(n\_farms=10, seed=42)}
    \item Solver: Gurobi QUBO with 30-second time limit per trial
    \item Lagrange multipliers tested: $\lambda \in \{1.0, 5.0, 10.0, 25.0, 50.0, 100.0, 150.0\}$
\end{itemize}

\textbf{Evaluation Metrics:}
\begin{itemize}
    \item \textbf{Constraint Violations:} Number of violated constraints (target: 0)
    \item \textbf{Objective Value:} Normalized weighted agricultural value
    \item \textbf{Land Utilization:} Percentage of total land assigned to crops
    \item \textbf{Crop Diversity:} Number of distinct crops planted
\end{itemize}

\subsubsection{Implementation Details}

The test script (\texttt{test\_lagrange\_multipliers.py}) performs the following workflow:

\begin{algorithm}[H]
\caption{Lagrange Multiplier Sensitivity Analysis}
\begin{algorithmic}[1]
\State Load food data and create scenario configuration
\State Generate 10-patch problem instance with fixed seed
\State Create CQM for binary formulation
\For{$\lambda \in \{\text{multipliers}\}$}
    \State Convert CQM to BQM using $\lambda$ as penalty weight
    \State Solve BQM with Gurobi QUBO (30s limit)
    \State Validate constraints and calculate metrics
    \State Record: violations, objective, utilization, diversity
\EndFor
\State \textbf{Output:} Table showing trade-offs across $\lambda$ values
\State \textbf{Recommendation:} Smallest $\lambda$ achieving 0 violations
\end{algorithmic}
\end{algorithm}

\subsubsection{Theoretical Considerations}

\paragraph{Penalty Method Theory:} For a constrained problem:
$$\min f(\mathbf{x}) \quad \text{s.t.} \quad g_i(\mathbf{x}) \leq 0, \, i=1,\ldots,m$$

the penalized formulation is:
$$\min_{\mathbf{x}} \, f(\mathbf{x}) + \lambda \sum_{i=1}^m \max(0, g_i(\mathbf{x}))^2$$

\textbf{Convergence Property:} As $\lambda \to \infty$, the penalized solution converges to the constrained optimum, but:
\begin{itemize}
    \item Objective landscape becomes increasingly steep near constraint boundaries
    \item Numerical conditioning deteriorates (ill-conditioned Hessian)
    \item Solver performance degrades due to extreme coefficient ratios
\end{itemize}

\paragraph{Practical Selection Rule:} Choose the smallest $\lambda$ such that:
$$\max_{i=1,\ldots,m} \max(0, g_i(\mathbf{x}^*)) < \epsilon$$
where $\mathbf{x}^*$ is the solution and $\epsilon$ is a tolerance (typically $10^{-6}$).

\subsubsection{Experimental Results}

\textbf{Test Configuration:}
\begin{itemize}
    \item Problem instance: 10 patches, full food dataset
    \item Total land: 621.58 ha (from farm\_sampler)
    \item Solver: Gurobi QUBO with 30-second time limit per trial
    \item BQM size: 434 binary variables, 16,069 non-zero QUBO terms
    \item Formulation: Binary (even grid) with $Y\_{p,c}$ variables only
\end{itemize}

\textbf{Results:}

\begin{center}
\begin{tabular}{ccccc}
\hline
$\lambda$ & Violations & Objective & Utilization & Crops \\
\hline
1.0 & 10 & 0.5387 & 200.0\% & 20 \\
5.0 & 10 & 0.5388 & 200.0\% & 20 \\
10.0 & 10 & 0.5336 & 200.0\% & 20 \\
25.0 & 10 & 0.5600 & 200.0\% & 16 \\
50.0 & 10 & 0.5589 & 200.0\% & 16 \\
100.0 & 10 & 0.5399 & 200.0\% & 20 \\
150.0 & 10 & 0.5143 & 200.0\% & 20 \\
\hline
\end{tabular}
\end{center}

\textbf{Critical Finding:}

The experimental results reveal that \textbf{all tested Lagrange multipliers} $\lambda \in [1.0, 150.0]$ produced constraint violations in the binary formulation. Specifically:

\begin{itemize}
    \item \textbf{Violation pattern:} All 10 plots assigned to multiple crops simultaneously
    \item \textbf{Constraint:} "At most one crop per plot" ($\sum_c Y_{p,c} \leq 1$)
    \item \textbf{Utilization:} 200\% (double-assignment of all land)
    \item \textbf{Implication:} Binary formulation requires significantly higher penalty weights than continuous formulation
\end{itemize}

\textbf{Analysis and Interpretation:}

\begin{enumerate}
    \item \textbf{Formulation-Dependent Scaling:}
    \begin{itemize}
        \item Binary formulation has different constraint structure than continuous
        \item "At most one per plot" constraint is \textit{harder} to enforce via penalties
        \item Requires $\lambda \gg 150$ to achieve feasibility
    \end{itemize}
    
    \item \textbf{Penalty Weight Requirements:}
    \begin{itemize}
        \item Continuous formulation (PATCH runner): $\lambda = 5$-$10$ sufficient
        \item Binary formulation: $\lambda = 500$-$1000$ estimated (extrapolation)
        \item Factor of $\approx 50$-$100\times$ difference between formulations
    \end{itemize}
    
    \item \textbf{Solver Behavior at Low $\lambda$:}
    \begin{itemize}
        \item Gurobi exploits weak penalties to maximize objective
        \item Double-assigns plots to achieve higher raw objective value
        \item BQM energy includes penalties, but solver optimizes BQM (not original CQM)
    \end{itemize}
\end{enumerate}

\textbf{Revised Recommendations:}

Based on these findings, for the \textbf{binary (even grid) formulation}:

\begin{itemize}
    \item \textbf{Minimum Feasible (estimated):} $\lambda \geq 500$ (requires further testing)
    \item \textbf{Recommended Default:} $\lambda = 1000$ (conservative, ensures feasibility)
    \item \textbf{Alternative Approach:} Use native CQM solver (LeapHybridCQMSampler) which handles constraints explicitly without penalty method
    \item \textbf{For CQM→BQM conversion:} Increase $\lambda$ iteratively until zero violations achieved
\end{itemize}

\textbf{Practical Implications:}

\begin{enumerate}
    \item \textbf{Penalty Method Limitations:} This experiment demonstrates that penalty-based constraint handling is \textit{formulation-sensitive}. The binary formulation's discrete structure makes constraints harder to enforce via soft penalties.
    
    \item \textbf{Solver Choice Matters:} For binary formulations with complex constraints:
    \begin{itemize}
        \item \textbf{Preferred:} D-Wave CQM Sampler (explicit constraints, no penalty tuning)
        \item \textbf{Alternative:} Classical MILP solvers (PuLP/Gurobi with hard constraints)
        \item \textbf{Use with caution:} BQM conversion with penalty method (requires careful $\lambda$ tuning)
    \end{itemize}
    
    \item \textbf{Benchmark Implementation:} Based on this finding, the benchmark uses:
    \begin{itemize}
        \item \textbf{Primary method:} PuLP with Gurobi (hard constraints, no $\lambda$ needed)
        \item \textbf{Quantum-classical:} D-Wave CQM Sampler (explicit constraints)
        \item \textbf{BQM experiments:} Use $\lambda = 1000$ as starting point, validate results
    \end{itemize}
\end{enumerate}



\newpage
\subsection{Grid Refinement Analysis}\label{sub:refinement}

\subsubsection{Motivation and Research Question}

The binary (PATCH) formulation discretizes continuous land into fixed-size plots with binary assignment variables. This introduces an \textbf{approximation error} compared to the continuous formulation. The fundamental question is:

\begin{center}
\textit{How does grid refinement (number of plots) affect solution quality?}
\end{center}

\textbf{Hypothesis:} Finer grids ($n \to \infty$) should converge to the continuous optimum, but at the cost of increased problem size and solve time.

\subsubsection{Experimental Design}

\textbf{Grid Refinement Levels Tested:}
$$n \in \{5, 10, 25, 50, 100\}$$

\textbf{Comparison Framework:}

For each refinement level $n$:
\begin{enumerate}
    \item \textbf{Continuous Baseline:} Solve with $n$ farms using \textbf{uneven distribution} (realistic sizes from \texttt{farm\_sampler})
    \item \textbf{Discretized:} Solve with $n$ patches using \textbf{even grid} (equal-sized plots via \texttt{patch\_sampler})
\end{enumerate}

Both scenarios use the same total land area for fair comparison.

\textbf{Evaluation Metrics:}
\begin{itemize}
    \item \textbf{Objective Value:} $Z_{\text{cont}}$ (continuous) vs. $Z_{\text{disc}}$ (discretized)
    \item \textbf{Optimality Gap:} $\Delta = \frac{Z_{\text{cont}} -- Z_{\text{disc}}}{Z_{\text{cont}}} \times 100\%$
    \item \textbf{Solve Time:} Wall-clock time for PuLP/Gurobi
    \item \textbf{Time Ratio:} $t_{\text{disc}} / t_{\text{cont}}$
\end{itemize}

\subsubsection{Implementation Details}

The test script (\texttt{Grid\_Refinement.py}) executes the following workflow:

\begin{algorithm}[H]
\caption{Grid Refinement Convergence Study}
\begin{algorithmic}[1]
\State \textbf{Input:} Total land area $A_{\text{total}}$
\For{$n \in \{5, 10, 25, 50, 100, 200\}$}
    \State \textbf{// Continuous Baseline}
    \State Generate $n$ farms with uneven distribution (total area $A_{\text{total}}$)
    \State Load food data (2 foods per group for tractability)
    \State Create continuous CQM with $A_{f,c}$ and $Y_{f,c}$ variables
    \State Solve with PuLP/Gurobi $\to$ $Z_{\text{cont}}, t_{\text{cont}}$
    \State
    \State \textbf{// Discretized Formulation}
    \State Generate $n$ patches with even grid (total area $A_{\text{total}}$, equal plot size $a_p = A_{\text{total}}/n$)
    \State Load same food data
    \State Create binary CQM with $Y_{p,c}$ variables only
    \State Solve with PuLP/Gurobi $\to$ $Z_{\text{disc}}, t_{\text{disc}}$
    \State
    \State Compute gap: $\Delta = (Z_{\text{cont}} -- Z_{\text{disc}})/Z_{\text{cont}} \times 100\%$
    \State Record: $n, Z_{\text{cont}}, Z_{\text{disc}}, \Delta, t_{\text{cont}}, t_{\text{disc}}, t_{\text{disc}}/t_{\text{cont}}$
\EndFor
\State \textbf{Output:} Table and convergence analysis
\end{algorithmic}
\end{algorithm}

\subsubsection{Theoretical Analysis}

\paragraph{Approximation Error Bound:}

For a linear objective $f(A) = \sum_c B_c A_c$, the discretization error is:

$$\epsilon(n) = \left| \sum_c B_c A_c^{\text{opt}} -- \sum_c B_c \left(\sum_p a_p Y_{p,c}^{\text{opt}}\right) \right|$$

where $A_c^{\text{opt}}$ is the continuous optimal area and $Y_{p,c}^{\text{opt}}$ is the discrete solution.

\textbf{Upper Bound:} If minimum planting areas $A_{\min,c}$ dominate, the error is bounded by:
$$\epsilon(n) \leq \sum_c B_c \cdot a_p = O\left(\frac{A_{\text{total}}}{n}\right)$$

\textbf{Convergence Rate:} $\epsilon(n) = O(n^{-1})$.

\paragraph{Computational Complexity Trade-off:}\\

\textbf{Continuous Formulation:}
\begin{itemize}
    \item Variables: $2nc$ ($n$ farms, $c$ crops)
    \item Constraints: $O(nc)$
    \item Solve time: $O(2^{nc} \cdot \text{poly}(nc))$ (MILP worst-case)
\end{itemize}

\textbf{Binary Formulation:}
\begin{itemize}
    \item Variables: $nc$ (binary only)
    \item Constraints: $O(n + c)$ (fewer due to simpler structure)
    \item Solve time: $O(2^{nc} \cdot \text{poly}(nc))$ (BIP worst-case, but better LP relaxation)
\end{itemize}

\textbf{Expected Behavior:}
\begin{itemize}
    \item Small $n$ ($n \leq 25$): Binary faster (fewer variables, simpler constraints)
    \item Large $n$ ($n \geq 100$): Continuous may be faster (continuous relaxation tighter than binary LP relaxation)
\end{itemize}

\subsubsection{Experimental Results}

\textbf{Test Configuration:}
\begin{itemize}
    \item Total land: 100 ha (fixed for all refinement levels)
    \item Food dataset: 27 foods across 5 food groups
    \item Solver: PuLP with Gurobi backend
    \item Continuous: Uneven farm distribution (farm\_sampler)
    \item Discretized: Even grid with equal plot sizes (patch\_sampler)
\end{itemize}

\textbf{Convergence Results:}

\begin{center}
\begin{tabular}{cccccc}
\hline
$n$ & $Z_{\text{cont}}$ & $Z_{\text{disc}}$ & Gap (\%) & $t_{\text{cont}}$ (s) & $t_{\text{disc}}$ (s) \\
\hline
5 & 0.2590 & 0.2263 & 12.63 & 0.062 & 0.005 \\
10 & 0.2589 & 0.2427 & 6.27 & 0.017 & 0.009 \\
25 & 0.2587 & 0.2525 & 2.38 & 0.051 & 0.016 \\
50 & 0.2583 & 0.2558 & 0.95 & 0.089 & 0.034 \\
100 & 0.2575 & 0.2575 & 0.00 & 0.220 & 0.060 \\
\hline
\end{tabular}
\end{center}

\textbf{Key Observations:}

\begin{enumerate}
    \item \textbf{Convergence Validation:}
    \begin{itemize}
        \item Gap decreases monotonically: 12.63\% → 6.27\% → 2.38\% → 0.95\% → 0.00\%
        \item \textbf{Convergence rate:} Approximately $O(n^{-1})$ as predicted theoretically
        \item At $n = 100$: Zero gap, perfect convergence to continuous optimum
    \end{itemize}
    
    
    \item \textbf{Computational Performance:}
    \begin{itemize}
        \item \textbf{Binary consistently faster:} Time ratio ranges 0.08x to 0.38x
        \item At $n = 100$: Binary is 3.7$\times$ faster (0.060s vs. 0.220s)
        \item \textbf{No crossover observed:} Contrary to expectation, binary remains faster even at $n = 100$
        \item Likely due to simpler constraint structure in binary formulation
    \end{itemize}
    
    \item \textbf{Scalability Limits:}
    \begin{itemize}
        \item $n = 200$: Continuous formulation became infeasible
        \item Possible causes: Too many small farms violating minimum area constraints
        \item Binary formulation avoids this issue through discrete plot assignment
    \end{itemize}
\end{enumerate}

\textbf{Convergence Analysis:}

Fitting the gap data to the theoretical model $\epsilon(n) = C/n$:

\begin{center}
\begin{tabular}{ccc}
\hline
$n$ & Observed Gap (\%) & Predicted Gap (\%) \\
\hline
5 & 12.63 & 12.63 (fitted) \\
10 & 6.27 & 6.32 \\
25 & 2.38 & 2.53 \\
50 & 0.95 & 1.26 \\
100 & 0.00 & 0.63 \\
\hline
\end{tabular}
\end{center}

The fitted constant $C \approx 63.2$ shows good agreement with observed data, confirming $O(n^{-1})$ convergence.








% \newpage
% \subsection{Benders Decomposition:}

% We solve the mixed-integer problem by partitioning the binary selection variables $\mathbf{Y}$ (the master) from the continuous allocation variables $\mathbf{A}$ (the subproblem). The original compact problem is:

% \begin{align}
% \min_{A,Y}\quad & -Z(A) + \mathbf{d}^T Y
%     && \text{(equivalently maximize $Z$)} \\
% \text{s.t.}\quad
% & \sum_{c} A_{f,c} \le L_f
%     && \forall f \\
% & A_{f,c} \ge A^{\min}_c\, Y_{f,c}, \quad
%   A_{f,c} \le L_f\, Y_{f,c}
%     && \forall f,c \\
% & FG^{\min}_g \le \sum_{f}\sum_{c\in G_g} Y_{f,c} \le FG^{\max}_g
%     && \forall g \\
% & A_{f,c} \ge 0,\quad Y_{f,c}\in\{0,1\}.
% \end{align}


% For a fixed candidate $\bar Y$ the subproblem is an LP in $\mathbf{A}$:

% \begin{align}
% \phi(\bar Y) = \min_{A \ge 0}\; & -Z(A) \\
% \text{s.t.}\quad
% & \sum_{c} A_{f,c} \le L_f && \forall f, \\
% & A_{f,c} \ge A^{\min}_c\,\bar Y_{f,c}, \quad
%   A_{f,c} \le L_f\,\bar Y_{f,c} && \forall f,c.
% \end{align}


% Dual information from the subproblem produces feasibility cuts when the subproblem is infeasible and optimality cuts when it is feasible; these cuts are added to the master, which can be written compactly as:

% \begin{equation}
% \begin{aligned}
% \min_{Y \in {0,1}^{\cdot},;\theta} \quad
% & \mathbf{d}^T Y + \theta \\
% \text{s.t.}\quad
% & \text{(food-group and other combinatorial constraints on $Y$)} \\
% & \theta \ge \beta^k + \sum_{f,c} \alpha^k_{f,c} Y_{f,c}
% \quad \text{for each optimality cut }k, \\
% & \text{feasibility cuts (when generated).}
% \end{aligned}
% \end{equation}

% The practical iterative procedure is given in Algorithm \ref{alg:benders} below.
% \begin{algorithm}[H]
% \caption{Benders Decomposition Algorithm}
% \label{alg:benders}
% \begin{algorithmic}[1]
% \State \textbf{Initialize:} Master problem with no cuts; iteration counter $k \gets 0$.
% \Repeat
% \State Solve the master MILP to obtain candidate binary vector $\bar{Y}^{(k)}$ and master estimate $\theta^{(k)}$.
% \State Solve the subproblem LP for fixed $\bar{Y}^{(k)}$:
% \If{subproblem is \textbf{infeasible}}
% \State Obtain an infeasibility ray and derive a feasibility cut.
% \State Add the cut to the master and go to Step 2.
% \Else
% \State Subproblem is feasible; compute $\phi(\bar{Y}^{(k)})$ and extract corresponding dual variables.
% \State Construct an optimality cut using the duals and add it to the master.
% \EndIf
% \State Evaluate the optimality gap: $|\theta^{(k)} -- \phi(\bar{Y}^{(k)})|$.
% \If{gap $\le$ tolerance}
% \State \textbf{Stop:} Current solution is optimal.
% \Else
% \State $k \gets k + 1$.
% \EndIf
% \Until{convergence}
% \end{algorithmic}
% \end{algorithm}



% \begin{table}[ht]
% \centering
% \caption{Notation and symbols}
% \label{tab:notation}
% \small
% \begin{tabular}{@{} l p{8.5cm} l @{}} 
% \textbf{Symbol} & \textbf{Meaning} & \textbf{Domain / notes} \\ 
% $A_{f,c}$ & Area (continuous) allocated on farm $f$ to crop $c$. & $A_{f,c}\ge0$ (continuous) \\
% $Y_{f,c}$ & Binary selection: 1 if crop $c$ is planted on farm $f$, 0 otherwise. & $Y_{f,c}\in\{0,1\}$ \\
% $\mathbf{A},\:\mathbf{Y}$ & Stacked vectors/matrices of $A_{f,c}$ and $Y_{f,c}$ respectively. & -- \\
% $Z(A)$ & Continuous part of the objective (e.g., profit or yield) as a function of $A$. & sign convention: you minimize $-Z(A)$ in the model \\ 
% $\mathbf{d}$ & Vector of coefficients multiplying $Y$ in the objective (fixed costs or penalties). & $\mathbf{d}^T\mathbf{Y}$ appears in master objective \\ 
% $L_f$ & Total land (available area) on farm $f$. & scalar, farm-local upper bound \\ 
% $A^{\min}_c$ & Minimum area required for crop $c$ if selected on a farm. & used in coupling: $A_{f,c}\ge A^{\min}_c Y_{f,c}$ \\ 
% $G_g$ & Set of crops that belong to food-group $g$. & index set \\ 
% $FG^{\min}_g,\; FG^{\max}_g$ & Lower and upper bounds on the number of plantings (or coverage) for food-group $g$. & Integers or counts over $f,c\in G_g$ \\ 
% $\bar Y$ (or $\bar Y^{(k)}$) & Candidate binary vector returned by the master at an iteration (trial solution). & fixed when solving subproblem \\ 
% $\phi(\bar Y)$ & Optimal value of the subproblem (LP) for fixed $\bar Y$. & finite if subproblem feasible; otherwise indicates infeasibility \\ 
% $\theta$ & Auxiliary master variable that lower-bounds the subproblem value ($\theta\ge\phi(\cdot)$ via cuts). & continuous scalar in master \\ 
% $\alpha^k_{f,c}$ & Coefficients of $Y_{f,c}$ in optimality cut $k$ (derived from dual solution). & appear in $\theta \ge \beta^k + \sum_{f,c}\alpha^k_{f,c}Y_{f,c}$ \\ 
% $\beta^k$ & Constant (intercept) term in optimality cut $k$ (from duals and constants). & constant offset in cut $k$ \\ 
% $\pi_f$ & Dual multiplier for farm-level land constraint $\sum_c A_{f,c}\le L_f$. & $\pi_f\ge0$ (depending on convention) \\ 
% $f$ & Index for farms (element of index set $\mathcal{F}$). & discrete index \\
% $c$ & Index for crops (element of index set $\mathcal{C}$). & discrete index \\
% $g$ & Index for food-groups (element of index set $\mathcal{G}$). & discrete index \\
% $k$ & Index for Benders cuts / iterations (cut counter). & integer iteration / cut index \\
% infeasibility ray & Dual ray produced when the subproblem is infeasible; used to form feasibility cuts. & yields linear inequality in $Y$ \\ 
% \end{tabular}
% \end{table}


\newpage

\section*{Team Presentation}

\begin{center}
\begin{tabular}{ | m{5em} | m{4em}| m{5em} | m{15em} | m{8em} | } 
 \hline
 Team Member 
 (First name, Last name)
 & Affiliation & Country (of the affiliation) & Relevant domain expertise for the project 
(i) Quantum computing, 
ii) SDG domain, 
iii) Application domain, 
iv) Classical computation (e.g. AI, ML, chemistry, operation research, fluid dynamics, etc.), 
v) other)
& Short Bio (3-5 sentences) 
 \\ 
 \hline 
  &  &  &  & \\ 
 \hline
 &  &  &  & \\ 
 \hline
  &  &  &  & \\ 
 \hline
  &  &  &  & \\ 
 \hline
\end{tabular}
\end{center}

\bibliographystyle{alpha}
\bibliography{references}



\end{document}